% !TeX program = xelatex
% ======================================================================
% 2026 中国研究生数学建模竞赛论文
% 基于 format.cls 官方规范排版
% ======================================================================

\PassOptionsToPackage{quiet}{xeCJK}
\documentclass[withoutpreface,bwprint]{format}

% 算法宏包与中文配置
\usepackage{algorithm}
\usepackage{algpseudocode}
\floatname{algorithm}{算法}
\renewcommand{\algorithmicrequire}{\textbf{输入:}}
\renewcommand{\algorithmicensure}{\textbf{输出:}}

% 论文标题与关键词
\newcommand{\PaperTitle}{山区洪涝灾害下无人机应急物资运输与通信协同优化建模研究}
\newcommand{\PaperKeywords}{应急物资运输\quad 无人机组批\quad 非线性能耗\quad 最大安全载荷\quad 集合划分模型\quad 状态压缩动态规划\quad 敏感性分析}

% 图片统一从 figures 文件夹读取
\graphicspath{{figures/}{./}}
\DeclareGraphicsExtensions{.png,.jpg,.jpeg,.pdf}

% 表格列自动换行与居中
\newcolumntype{L}{>{\raggedright\arraybackslash}X}
\newcolumntype{C}{>{\centering\arraybackslash}X}

% 安全图片加载宏：优先查找 figures/，防止因缺图中断编译
\newcommand{\SafeGraphic}[2][width=0.8\textwidth]{%
    \IfFileExists{figures/#2}{%
        \includegraphics[#1]{figures/#2}%
    }{%
        \IfFileExists{#2}{%
            \includegraphics[#1]{#2}%
        }{%
            \begin{center}
            \fbox{\parbox[c][4cm][c]{0.85\textwidth}{%
                \centering \song\color{gray}
                \textbf{【图占位符：\detokenize{#2}】}\\[0.5em]
                \small （提示：请在 figures 目录下放置图片文件 \texttt{\detokenize{#2}}）%
            }}
            \end{center}
        }%
    }%
}

% 取消 format.cls 自带的自动章节编号前缀，保留原稿中的大纲标题编号（如“一、”、“1.1”、“5.1.1”）
\makeatletter
\def\@seccntformat#1{}
\makeatother
\renewcommand{\thesection}{}
\renewcommand{\thesubsection}{}
\renewcommand{\thesubsubsection}{}

\begin{document}
\song\zihao{-4}

% 论文题目：黑体三号、居中
{\centering\zihao{3}\bfseries\heiti \PaperTitle\par}
\setcounter{page}{1}
\thispagestyle{empty}

% 摘要与关键词
\begin{abstract}
2026 年 7 月，台风“美莎克”带来的持续强降雨使广西横州市镇龙乡出现断路、断电、断网，约 8000 人一度被困。地面车辆无法进入，物资投送只能依靠无人机。本文以该场景为背景，建立了“航路高程与能耗模型—单架次可行能力—货箱组批优化—多目标权衡与敏感性分析”的求解链条，并给出问题一的完整结果。

问题一不考虑实体无人机与共享电池的跨架次调度，只研究 $O01 \to S_i \to O01$ 的单点往返。本文做了四方面工作。

\begin{enumerate}[leftmargin=*]
    \item \textbf{航路与能耗模型}。由 30 m 分辨率 DEM 沿大圆航线等间距采样，取沿线最高地面高程加 50 m 作为计划巡航海拔；能耗按水平巡航能耗与爬升附加能耗两部分计算，其中水平部分依题面 $3/2$ 次方航程衰减律标定（题面未给出水平能耗的显式公式，故按等效单位距离电耗标定并作 $\pm 20\%$ 敏感性检验），爬升部分按重力势能折算。

    \item \textbf{最大安全载荷}。将额定载重、航程校核与返航安全余量（基准 $\rho=20\%$）写成一个隐式非线性约束，用 60 次二分迭代求解 $15$ 个服务区 $\times$ $3$ 种机型的最大安全载荷矩阵。结果为：A 型在全部 15 个服务区都达到结构上限 $25\,\text{kg}$；B 型仅在最远的 S008 因电量不足降到 $28.63\,\text{kg}$；C 型在 5 个距离超过 $7.2\,\text{km}$ 的服务区受电量约束，载荷降至 $58.44\sim68.55\,\text{kg}$。可见本场景中限制投送能力的主要是电池能量和结构载重，航程校核始终没有起作用。

    \item \textbf{货箱组批}。以服务区为独立单元建立集合划分模型，约束为质量、体积和返航电量三条。由于单个服务区最多 15 箱，用位掩码状态压缩动态规划即可精确求解，无需启发式算法。

    \item \textbf{多目标权衡与敏感性}。比较“架次优先”“时间优先”“能耗优先”三种字典序策略。基准下前两者给出 18 架次、$59.240\,\text{kWh}$、$32798.5\,\text{s}$ 的方案，“能耗优先”给出 19 架次、$59.142\,\text{kWh}$、$34466.5\,\text{s}$ 的方案，两者互不支配：多飞一架次省下 $0.098\,\text{kWh}$，却多花 $1668\,\text{s}$ 地面作业时间。综合时效性，本文推荐 18 架次方案，80 个货箱全部交付，各架次返航 SOC 在 $22.79\%\sim73.07\%$ 之间，均不低于 $20\%$。敏感性方面，返航余量在 $10\%\sim20\%$ 区间内组批方案不变；升到 $25\%$ 时 S004 的一架次被迫拆成两架次，总架次增至 19；升到 $30\%$ 时增至 20。
\end{enumerate}

问题二解除单点往返限制，允许一个架次串联访问多个服务区，并引入 8 架实体机、14 组共享电池的两阶段充电周转与物资时限约束。本文做了三方面工作。

\begin{enumerate}[leftmargin=*]
    \item \textbf{路线模式与排程建模}。把“哪些箱同批、按什么顺序访问、用哪型飞机”离线固化成路线模式，模式内部用 6.2 节的物理模型一次算清逐箱交付偏移、时长、能耗与充电时间；模型里只保留“选哪些模式、各自何时开始”两类决策，配以货箱恰好覆盖一次、46 项硬时限、各型实体机与电池的并发容量、完工时间五组约束。同型资源完全可互换，因此只需约束并发占用数，具体编号在求解后用区间着色补出。

    \item \textbf{字典序求解与反馈迭代}。按完工时间、加权迟到量、总能耗、架次数四级依次优化，每级求完即把取值冻结为上界，保证后一级不牺牲前一级。模式库先由三种逐区精确划分、构造式组批、两区合并和全部单箱模式合成，保留所有物理可行机型、不做能耗支配剪枝；再按“拆分航线、单箱转移”的邻域逐轮扩充，两轮后模式库由 337 条增至 457 条。

    \item \textbf{最终方案}。28 架次方案（15A+7B+6C，含 2 条两点回路）把完工时间压到 $6347.2\,\text{s}$，总能耗 $71.970\,\text{kWh}$；80 件货箱全部交付，46 项硬时限全部达标，加权迟到量为 $0$，即连软时限也全部按期送达。全部 58 个飞行航段（32 条有向航段）经独立几何审计，最小净空 $50.000\,\text{m}$，违规数为 0。方案结构上值得注意的一点是，架次虽比问题一的 18 架次多，完工时间反而更短——本问真正的瓶颈是资源周转而不是飞行时间。
\end{enumerate}

本文同时搭建了后续问题三（通信中继协同）与问题四（任务空间分区与资源配置）的建模接口。
\end{abstract}
\keywords{\PaperKeywords}
\thispagestyle{empty}

% 正文
\section{一、问题重述}

\subsection{1.1 现实背景}
2026 年 7 月，受台风“美莎克”外围环流影响，广西横州市镇龙乡出现持续强降雨，引发山洪与滑坡。据公开报道，该乡一度出现断路、断电、断网，全乡约 8000 人被困，大型无人机承担了电力抢修材料运输，临时通信设备用于逐步恢复联络。这一场景的典型特征是：地面交通与地面通信同时失效，物资投送和通信保障都必须在道路中断、能量受限的条件下组织。灾后初期的救援需求往往同时具备突发、时限紧、基础设施损毁严重等特点，属于典型的人道主义物流问题\cite{caunhye2012, holguin2012}；而无人机受地面道路条件约束小，可以在道路中断时承担投送任务\cite{murray2015, otto2018}。

与常规物流相比，这类救援任务给建模带来三个具体困难。

\begin{enumerate}[leftmargin=*]
    \item \textbf{地形与能耗耦合}。山区高程起伏大，航线必须避开沿线山体并保持净空；载荷又会抬高多旋翼的悬停与巡航功率，电池能量非常有限，载荷与航程之间存在非线性关系。
    \item \textbf{货箱离散}。物资以不可拆分的货箱为单位，必须在载质量、货舱体积和返航电量三重边界内完成分配，既不能拆分也不能跨服务区合并。
    \item \textbf{多目标冲突}。现场既希望少出架次（减少空域协调与起降调度压力），又希望总电耗低（前线充电能力有限），还希望作业时间短。三者往往不能同时最优。
\end{enumerate}

\subsection{1.2 标准测试场景与数据}
赛题以镇龙乡及周边山区为背景，给出了空间、装备和物资三类数据。

\begin{enumerate}[leftmargin=*]
    \item \textbf{节点数据}：1 个临时调度中心 $O01$ 和 15 个受灾服务区 $S001\sim S015$ 的经纬度与地面海拔，以及覆盖全区的 30 m 分辨率 DEM。
    \item \textbf{装备数据}：A、B、C 三种运输无人机，在空载/满载标准航程、电池可用能量、最大载货质量、货舱容积、巡航与爬升/下降速度、返航电量下限、地面作业时间参数等方面互不相同。
    \item \textbf{物资数据}：医疗物资、饮用水、应急食品、生活卫生用品四类共 80 个货箱，逐箱给出所属服务区、单箱质量和单箱体积。
\end{enumerate}

\subsection{1.3 各问题的逻辑关系}
四个问题依次递进，共享同一套物理规则和数据口径。

\begin{enumerate}[leftmargin=*]
    \item \textbf{问题一}：不考虑实体无人机排班与电池周转，研究 $O01 \to S_i \to O01$ 的单点往返能力与货箱组批，是后三问的物理基础。
    \item \textbf{问题二}：解除单点往返限制，允许一个架次串联访问多个服务区，并加入实体机数量、共享电池周转和两阶段充电约束。把“箱组 + 访问序 + 机型”离线固化为路线模式，再用约束规划在同一模式库上联合选择模式与开始时刻，求解多目标运输调度。
    \item \textbf{问题三}：在问题二基础上加入通信链路预算与 DEM 视线遮挡判定，当直连被山体阻断时调度中继无人机，实现运输与通信的联合排程。
    \item \textbf{问题四}：以问题三的方案为输入，把 15 个服务区划分为 $K=2$ 和 $K=3$ 个任务组，按峰值并发核算各执行单元所需的机体、电池与中继资源。
\end{enumerate}



\section{二、模型基本假设}

赛题附录 2 已经规定了航段划分、等效航程公式、飞行时间构成、能耗构成与返航安全余量，这些直接实现，不再列入假设。以下 7 条是题面未明确、需要本文声明的内容。

\begin{enumerate}[label=\textbf{假设 \arabic*：}, leftmargin=*]
    \item \textbf{水平巡航能耗的标定方式。} 题面给出了载荷—等效航程的 $3/2$ 次方关系和电池可用能量 $E_g^{use}$，但没有给出水平巡航能耗的显式公式。本文按“可用能量除以对应等效航程”的方式标定单位距离电耗，即 $\varepsilon_g(q)=E_g^{use}/L_g(q)$；等价地，空载单位距离电耗为 $\varepsilon_g^0=E_g^{use}/L_g^0$，载重 $q$ 时按 $L_g^0/L_g(q)$ 的比例放大。这一标定与旋翼诱导功率随全重呈 $P_{ind}\propto M^{3/2}$ 变化的动量理论结论\cite{leishman2006}方向一致，也与 Dorling 等\cite{dorling2017}、Zhang 等\cite{zhang2021}、Stolaroff 等\cite{stolaroff2018}关于载荷显著抬高单位距离电耗的结论相符。需要说明的是，它是一个工程等效标定，不是唯一的空气动力学模型，因此本文另设 $\pm20\%$ 的水平能耗放缩做敏感性检验。

    \item \textbf{航程校核。} 题面规定 $L_g(q)$ 是等效航程，但未说明它如何约束一个往返架次。本文除总能耗约束外，额外要求去程单程水平距离不超过去程载荷对应的等效航程 $L_g(q)$，返程不超过空载标准航程 $L_g^0$。数值结果表明这一校核在本场景中从未成为紧约束（最紧的 S008 A 型仍余 $11.9\,\text{km}$），保留它只是为了避免在更极端参数下出现物理上不可达的方案。

    \item \textbf{爬升附加能耗的口径。} 题面给出了参数“爬升能耗效率”$=0.72$，但未给出爬升能耗的公式。本文把它解释为无量纲的推进效率，爬升附加能耗按重力势能折算：$E^{up}=mgh^+/(3.6\times10^6\eta)$，其中 $m$ 为当前总质量，$3.6\times10^6$ 是焦耳到 kWh 的换算常数。按此折算，爬升能耗在总能耗中占比约 $4\%$，即使该参数口径有出入，对结果影响也有限。下降阶段按题面规定取效率 0，不计附加能耗。

    \item \textbf{水平距离口径。} 题面称节点间采用“水平直线”，但节点坐标是经纬度。本文将其理解为球面大圆距离，地球半径取 $R=6\,371\,008.8\,\text{m}$\cite{sinnott1984}。本场景中服务区距 $O01$ 仅 2.8--8.1 km，大圆距离与椭球测地线\cite{vincenty1975}、平面近似之间的差异都在米级以下，不影响结论；之所以仍采用大圆公式，是为了与后续问题的三维通信链路计算保持同一几何口径。

    \item \textbf{累计作业时间的口径。} 题面只要求“累计作业时间”，未定义其构成。本文参考多旋翼作业时间的一般构成\cite{liu2017}，把每个架次的作业时间定义为：固定出动准备时间 $+$ 单箱装载时间 $\times$ 箱数 $+$ 往返飞行时间 $+$ 到场交接基础时间 $+$ 单箱交接时间 $\times$ 箱数，参数全部取自附件。这一口径下地面保障时间与架次数强相关，为便于读者判断，本文在 5.5.2 节同时报告纯飞行时间。

    \item \textbf{多目标的处理方式。} 题面要求综合三个指标但未给权重。本文不引入主观权重，改用字典序比较三种代表性策略（架次优先、能耗优先、时间优先），据此说明指标间的取舍。报告的是这三个代表解之间的非支配关系，不是完整 Pareto 前沿的证明。

    \item \textbf{高程与作业高度的取值。} 计划巡航海拔取航段所经 DEM 像元的最高地面高程加 $50\,\text{m}$；$O01$ 作业高度取节点地面海拔，服务区作业高度取地面海拔加 $30\,\text{m}$。DEM 按 EPSG:4326 处理。节点表自带海拔与 DEM 在该点的取值可能存在少量不一致，本文一律以节点表为准，未作调和。

    \item \textbf{实体机与共享电池初始就绪状态（针对问题二）。} 设前线临时调度中心 $O01$ 在初始 $t=0$ 时刻拥有全部 8 架实体运输无人机（U01$\sim$U04 为 A 型，U05$\sim$U06 为 B 型，U07$\sim$U08 为 C 型）和全部 14 组共享电池（A-B01$\sim$A-B06 为 A 型，B-B01$\sim$B-B04 为 B 型，C-B01$\sim$C-B04 为 C 型）。所有机体与电池均处于可用状态，初始电池荷电状态 $\text{SOC}=100\%$；基地具备充足的并发充电接口，不同电池的充电过程相互独立、互不干扰。

    \item \textbf{物资时限的硬/软约束划分与及时性度量（针对问题二）。} 将 80 件货箱的时限划分为两个层级：全部 16 件医疗物资的期望送达时间构成 16 项硬约束；标注为“首批保障”的货箱共 30 件（每个服务区 2 件，1 件医疗物资加 1 件饮用水），其首批截止时间构成 30 项硬约束，其中 15 件与前述医疗箱重合，故硬约束合计 $16+30=46$ 项，在排程中必须无条件满足，不允许任何延迟。其余 64 件非医疗货箱的期望送达时间作为服务质量“软指标”，以应急优先级系数加权迟到量进入次级优化目标，兼顾救援的刚性红线与整体调度的灵活性。

    \item \textbf{多架次调度的首要完成时间口径（针对问题二）。} 题目所指问题二优化目标中的“时间最短”，明确界定为\textbf{全部任务完成时间（Makespan $C_{\max}$）}，即机队所有架次完成全部投送任务并返回 $O01$ 的最晚时刻（$\max_{p \in \mathcal{P}} r_p$），而非各架次作业时间之和，亦非最晚货箱交付时刻。在保证最短完工时间前提下，再依次优化非医疗加权迟到量、总能耗与架次数。
\end{enumerate}



\section{三、符号与参数说明}

表 \ref{tab:notations} 汇总全文使用的主要符号。

\begin{table}[H]
\centering
\small
\caption{全文核心数学符号与参数定义}
\label{tab:notations}
\begin{tabularx}{\textwidth}{p{3.2cm} p{2.2cm} X}
\toprule
\textbf{符号} & \textbf{量纲/单位} & \textbf{物理定义与说明} \\
\midrule
\multicolumn{3}{l}{\textbf{集合与索引}} \\
$\mathcal{S} = \{S_1, \dots, S_{15}\}$ & -- & 受灾服务区集合，共 15 个目标点 \\
$\mathcal{G} = \{A, B, C\}$ & -- & 异构运输无人机机型集合 \\
$\mathcal{U} = \{U_1, \dots, U_8\}$ & -- & 实体运输无人机集合（4架A型、2架B型、2架C型） \\
$\mathcal{K} = \{B_1, \dots, B_{14}\}$ & -- & 共享动力电池集合（6组A型、4组B型、4组C型） \\
$\mathcal{B}$ & -- & 全部待运送应急货箱集合，$|\mathcal{B}| = 80$ \\
$\mathcal{B}_i$ & -- & 目的地为受灾服务区 $S_i$ 的货箱子集，$\bigcup \mathcal{B}_i = \mathcal{B}$ \\
$\mathcal{R}_i$ & -- & 服务区 $S_i$ 的合法可行候选架次集合 \\
$\mathcal{P}, p$ & -- & 全局协同调度执行架次集合与架次索引 \\
$i, j \in \mathcal{S} \cup \{O01\}$ & -- & 地理节点索引（$O01$ 为临时调度中心） \\
$g \in \mathcal{G}$ & -- & 运输无人机机型索引 \\
$b \in \mathcal{B}$ & -- & 货箱编号索引 \\
$r \in \mathcal{R}_i$ & -- & 候选架次索引 \\
\midrule
\multicolumn{3}{l}{\textbf{地理空间与航路参数}} \\
$d_{ij}$ & $\text{m}$ & 节点 $i$ 与节点 $j$ 之间的地表球面大圆水平距离 \\
$H_i^{ground}$ & $\text{m}$ & 节点 $i$ 的地面高程（海拔） \\
$H_{ij}^{max}$ & $\text{m}$ & 航线 $i \to j$ 正下方地表沿 DEM 剖面的最高地面海拔 \\
$H_{ij}^{cruise}$ & $\text{m}$ & 航线 $i \to j$ 的巡航计划绝对海拔，$H_{ij}^{cruise} = H_{ij}^{max} + 50$ \\
$h_{ij}^+, h_{ij}^-$ & $\text{m}$ & 航段 $i \to j$ 的垂直爬升高度与垂直下降高度 \\
\midrule
\multicolumn{3}{l}{\textbf{无人机装备物理属性参数}} \\
$m_g^{empty}$ & $\text{kg}$ & 机型 $g$ 的含电池空载总质量 \\
$Q_g$ & $\text{kg}$ & 机型 $g$ 的最大允许载货质量（结构载重上限） \\
$V_g$ & $\text{m}^3$ & 机型 $g$ 货舱的最大有效装载容积 \\
$v_g^c, v_g^\uparrow, v_g^\downarrow$ & $\text{m/s}$ & 机型 $g$ 的水平巡航速度、垂直爬升速度与垂直下降速度 \\
$L_g^0, L_g^F$ & $\text{m}$ & 机型 $g$ 在标称电池电量下的空载标准航程与满载标准航程 \\
$E_g^{use}$ & $\text{kWh}$ & 机型 $g$ 配套单组动力电池的额定可用能量容量 \\
$T_{full, g}$ & $\text{s}$ & 机型 $g$ 动力电池 0--100\% 等效完全充电基准时间 \\
$\rho_g$ & -- & 机型 $g$ 的返航安全预留电量比例（基准值取 $0.20$） \\
$\eta_g$ & -- & 机型 $g$ 垂直爬升推进能耗折算效率（标定为 $0.72$） \\
$\tau_g^{prep}, \tau_g^{load}$ & $\text{s}$ & 单架次出动固定准备时间、单箱物资装载耗时 \\
$\tau_g^{base}, \tau_g^{drop}$ & $\text{s}$ & 目的地现场交接基础准备耗时、单箱货物交接作业耗时 \\
$D_b^{hard}, D_b^{exp}$ & $\text{s}$ & 货箱 $b$ 的硬截止时间与一般期望送达时间 \\
$w_b$ & -- & 货箱 $b$ 的应急优先级加权系数 \\
\midrule
\multicolumn{3}{l}{\textbf{决策变量与状态评估函数}} \\
$Q_{safe, ig}$ & $\text{kg}$ & 机型 $g$ 针对服务区 $S_i$ 执行单点往返的最大安全载荷 \\
$q_r, v_r$ & $\text{kg}, \text{m}^3$ & 候选架次 $r$ 的实际承载货物总质量与总体积 \\
$E_{ir}, T_{ir}$ & $\text{kWh}, \text{s}$ & 候选架次 $r$ 执行单点往返的总能量消耗与累计作业时间 \\
$\text{SOC}_r^{ret}, \text{SOC}_p^{ret}$ & $\%$ & 架次返抵 $O01$ 时电池剩余荷电状态百分比 \\
$s_p, r_p$ & $\text{s}$ & 架次 $p$ 的准备开始时刻与返航着陆时刻 \\
$c_p$ & $\text{s}$ & 架次 $p$ 返航后将电池充至 100\% 所需的充电耗时 \\
$t_{pb}$ & $\text{s}$ & 货箱 $b$ 在架次 $p$ 中投送完成的绝对时刻 \\
$C_{\max}$ & $\text{s}$ & 机队全部架次返航的最晚完成时间（Makespan） \\
$x_{ir}, y_{pb} \in \{0, 1\}$ & -- & 0-1 决策变量：货箱分配与架次执行状态变量 \\
\bottomrule
\end{tabularx}
\end{table}



\section{四、地理空间环境与数据预处理}

\subsection{4.1 节点拓扑与高程栅格核验}
根据附件《调度中心与服务区.xlsx》，$O01$（凤丹村）位于东经 $109.2308517^\circ$、北纬 $23.0085095^\circ$，地面海拔 $127.7\,\text{m}$。15 个服务区地面海拔在 $154.0\,\text{m}$（S001）至 $444.5\,\text{m}$（S015）之间。到 $O01$ 的水平大圆距离最短为 S011 的 $2.82\,\text{km}$，最远为 S008 的 $8.11\,\text{km}$；15 个服务区合计需保障人口 3422 人，其中 S001 一个点就占 2100 人。

DEM 文件《镇龙乡及周边30米DEM.tif》读取结果为：投影类型 geographic，地理坐标系代码 4326（即 WGS84），像元尺度 $0.0002778^\circ$，在本地纬度上约合 $30\,\text{m}$（纬向 $30.9\,\text{m}$、经向 $28.5\,\text{m}$）；数据为 32 位浮点。栅格覆盖范围约在东经 $109.03^\circ\sim109.45^\circ$、北纬 $22.86^\circ\sim23.22^\circ$ 之间，$O01$ 与全部服务区均落在覆盖范围内。需要说明的是，该栅格没有标注 NoData 值，程序中对越界采样作了显式报错处理，但 NoData 校验在本数据上实际不会触发。

\subsection{4.2 航路三维剖面与高程采样}
在山区低空飞行中，沿线地形是决定爬升高度进而决定能耗的主要因素\cite{roberge2013}。按附录 2 的规则，航段为两点间的水平直线，巡航海拔取该航段所经 DEM 像元的最高地面高程加 $50\,\text{m}$。

设航段端点经纬度为 $P_1$、$P_2$，水平距离为 $d_{12}$。本文以不超过 $30\,\text{m}$ 的步长在两点间均匀布设采样点：

\begin{equation}
N_{step} = \max\left(1, \left\lceil \frac{d_{12}}{30} \right\rceil\right), \quad P(k) = P_1 + \frac{k}{N_{step}} (P_2 - P_1), \quad k \in \{0, 1, \dots, N_{step}\}
\end{equation}

每个采样点的经纬度先转换为栅格行列号，再用双线性内插取值，得到沿线最高地表高程：

\begin{equation}
H_{max} = \max_{k \in \{0, \dots, N_{step}\}} z(k)
\end{equation}

由此得到巡航海拔与两端爬升、下降高度：

\begin{equation}
H^{cruise} = H_{max} + 50\,\text{m}, \quad h^+ = \max(0, H^{cruise} - H_{start}), \quad h^- = \max(0, H^{cruise} - H_{end})
\end{equation}

去程起点为 $O01$，作业海拔 $127.7\,\text{m}$；终点为服务区，作业高度取地面海拔加 $30\,\text{m}$。返程时两端角色互换。

沿线地形对结果的影响并不总是随距离增加。以 S014 为例，该点自身地面高程为 $321\,\text{m}$，但航路沿线最高地面高程达 $541.95\,\text{m}$，高出起点地面 $220.95\,\text{m}$，对应巡航海拔 $591.95\,\text{m}$，是 15 条航路中最高的一条，其爬升高度甚至超过直线距离更远的 S008（巡航海拔 $372.91\,\text{m}$）。这说明在山区场景中，单看水平距离会明显低估能耗，必须结合地形剖面。

\subsection{4.3 货箱属性与完整性审计}
对《物资需求与配送时限.xlsx》的 80 个货箱逐条核对，箱号唯一、服务区归属完整，未发现缺失或重复。汇总如下。

\begin{enumerate}[leftmargin=*]
    \item \textbf{总量}：80 箱，总质量 $758.0\,\text{kg}$，总体积 $2.011\,\text{m}^3$。
    \item \textbf{物资类别}：医疗物资 16 箱（每箱 $3\,\text{kg}$、$0.012\,\text{m}^3$）；饮用水 36 箱（每箱 $14\,\text{kg}$、$0.027\,\text{m}^3$）；应急食品 19 箱（每箱 $8\,\text{kg}$、$0.028\,\text{m}^3$）；生活卫生用品 9 箱（每箱 $6\,\text{kg}$、$0.035\,\text{m}^3$）。同类货箱的规格完全一致。
    \item \textbf{空间分布}：需求分布很不均匀。S001 最多，15 箱、$154.0\,\text{kg}$、$0.394\,\text{m}^3$；S002 与 S003 各 8 箱、$81.0\,\text{kg}$；S004、S005、S006 各 6 箱、$59.0\,\text{kg}$；S007、S008 各 5 箱、$45.0\,\text{kg}$；S009 至 S015 各 3 箱、$25.0\,\text{kg}$、$0.067\,\text{m}^3$。
\end{enumerate}



\section{五、问题一：单点往返运输能力与货箱组批优化模型}

\subsection{5.1 航段飞行剖面与能耗模型}

\subsubsection{5.1.1 飞行时间}
一个航段的飞行由垂直爬升、水平巡航、垂直下降三段组成。机型 $g$ 在航段 $i \to j$ 上的飞行时间按附录 2 的公式计算：

\begin{equation}
t_{gij} = \frac{h_{ij}^+}{v_g^\uparrow} + \frac{d_{ij}}{v_g^c} + \frac{h_{ij}^-}{v_g^\downarrow}
\end{equation}

单点往返架次 $O01 \to S_i \to O01$ 的飞行时间为去程与返程之和：

\begin{equation}
t_g^{flight} = t_{g, O01 \to S_i} + t_{g, S_i \to O01}
\end{equation}

\subsubsection{\texorpdfstring{5.1.2 载荷—等效航程的 $3/2$ 次方关系}{5.1.2 载荷—等效航程的 3/2 次方关系}}
附录 2 给出的等效航程公式为：

\begin{equation}
L_g(q) = L_g^0 - (L_g^0 - L_g^F) \left(\frac{q}{Q_g}\right)^{3/2}, \quad \forall q \in [0, Q_g]
\label{eq:range_32}
\end{equation}

式中 $L_g^0$、$L_g^F$ 分别为空载与满载标准航程，$Q_g$ 为额定载重。该式表明载重增加时等效航程的衰减快于线性：以 C 型为例，载重达到额定值的 $56\%$（$45\,\text{kg}$）时，等效航程已从 $26\,\text{km}$ 降到 $20.1\,\text{km}$。这一非线性是后文“拆分反而省电”现象的来源之一。其物理背景是动量理论中旋翼诱导功率随全重呈 $P_{ind}\propto M^{3/2}$ 变化\cite{leishman2006}。

\subsubsection{5.1.3 水平巡航能耗的标定}
由于题面未给出水平巡航能耗的显式公式，本文按假设 1 用“可用能量 / 等效航程”标定载荷 $q$ 时的单位距离电耗：

\begin{equation}
\varepsilon_g(q) = \frac{E_g^{use}}{L_g(q)} \quad (\text{kWh/m})
\end{equation}

于是航段水平距离为 $d_{ij}$ 时的水平巡航能耗为：

\begin{equation}
E_{gij}^{hor}(q) = \varepsilon_g(q) \cdot d_{ij} = \frac{E_g^{use} \cdot d_{ij}}{L_g(q)}
\end{equation}

\subsubsection{5.1.4 爬升附加能耗与返航约束}
无人机由作业高度爬升到巡航高度需要额外做功。去程全重为 $m_g^{empty} + q$，返程货物已交付，全重恢复为 $m_g^{empty}$。按假设 3：

\begin{equation}
E_{gij}^{up}(q) = \frac{(m_g^{empty} + q) \cdot g \cdot h_{ij}^+}{3.6 \times 10^6 \cdot \eta_g} \quad (\text{kWh})
\end{equation}

其中 $g = 9.80665\,\text{m/s}^2$，$\eta_g = 0.72$。下降阶段按题面规定不计附加能耗。

于是机型 $g$ 执行一次单点往返（去程载重 $q$）的总电耗为：

\begin{align}
E_{ig}^T(q) &= E_{g, O01 \to S_i}(q) + E_{g, S_i \to O01}(0) \notag \\
&= \left[ \frac{E_g^{use} d_{O01, S_i}}{L_g(q)} + \frac{(m_g^{empty} + q) g h_{O01 \to S_i}^+}{3.6 \times 10^6 \cdot \eta_g} \right] + \left[ \frac{E_g^{use} d_{O01, S_i}}{L_g^0} + \frac{m_g^{empty} g h_{S_i \to O01}^+}{3.6 \times 10^6 \cdot \eta_g} \right]
\end{align}

架次总能耗须满足返航安全余量，即返抵 $O01$ 时剩余电量不低于 $\rho_g$：

\begin{equation}
E_{ig}^T(q) \le (1 - \rho_g) E_g^{use} \iff \text{SOC}_{return} = 1 - \frac{E_{ig}^T(q)}{E_g^{use}} \ge \rho_g = 20\%
\end{equation}

在本场景的基准方案中，水平巡航能耗合计 $56.685\,\text{kWh}$，占总能耗的 $95.7\%$；爬升附加能耗合计 $2.555\,\text{kWh}$，占 $4.3\%$。

\subsection{5.2 最大安全载荷}

\subsubsection{5.2.1 定义}
机型 $g$ 面向服务区 $S_i$ 的最大安全载荷，定义为同时满足结构载重、整架次返航能量和航程校核三项条件的最大连续载荷：

\begin{equation}
Q_{safe, ig} = \max \left\{ q \;\middle|\; 0 \le q \le Q_g, \; E_{ig}^T(q) \le (1 - \rho_g) E_g^{use}, \; d_{O01, S_i} \le L_g(q), \; d_{O01, S_i} \le L_g^0 \right\}
\end{equation}

\subsubsection{5.2.2 二分求解}
等效航程 $L_g(q)$ 关于 $q$ 单调递减，因此水平巡航能耗关于 $q$ 单调递增；爬升能耗也随 $q$ 线性递增。故总能耗关于 $q$ 单调递增，$Q_{safe,ig}$ 是唯一的临界截断点。求解时先检查两个端点：若空载即超能量上限，该机型对该点不可行；若满载仍满足约束，则直接取结构上限 $Q_g$；否则在 $[0, Q_g]$ 上二分。算法 \ref{alg:bisection} 取 60 次迭代，收敛误差上限为 $Q_g \cdot 2^{-60} < 10^{-16}\,\text{kg}$，已到双精度浮点的分辨率极限。

\begin{algorithm}[H]
\caption{基于能耗二分搜索的最大安全载荷求解}
\label{alg:bisection}
\begin{algorithmic}[1]
\Require 机型参数 $g$、调度中心 $O01$、服务区 $S_i$、航路地表最高高程 $H_{max}$、安全余量 $\rho_g$
\Ensure 最大安全有效载荷 $Q_{safe, ig}$ 及对应的飞行状态
\State 计算大圆水平距离 $d = d(O01, S_i)$ 与巡航海拔 $H^{cruise} = H_{max} + 50\,\text{m}$
\State 计算空载往返能耗 $E_{low} = E_{ig}^T(0)$
\If{$E_{low} > (1 - \rho_g) E_g^{use}$ \textbf{or} $d > L_g^0$}
    \State \Return $Q_{safe, ig} = 0.0$ \Comment{能量不足或航程超限}
\EndIf
\State $q_{low} \gets 0.0$, $q_{high} \gets Q_g$
\For{$step = 1 \to 60$}
    \State $q_{mid} \gets (q_{low} + q_{high}) / 2.0$
    \State 计算 $L_g(q_{mid})$ 与整架次总能耗 $E_{ig}^T(q_{mid})$
    \If{$E_{ig}^T(q_{mid}) \le (1 - \rho_g) E_g^{use}$ \textbf{and} $d \le L_g(q_{mid})$}
        \State $q_{low} \gets q_{mid}$ \Comment{可行，向上探测}
    \Else
        \State $q_{high} \gets q_{mid}$ \Comment{超标，向下收缩}
    \EndIf
\EndFor
\State \Return $Q_{safe, ig} = q_{low}$
\end{algorithmic}
\end{algorithm}

\subsubsection{\texorpdfstring{5.2.3 $15\times 3$ 最大安全载荷计算结果}{5.2.3 15x3 最大安全载荷计算结果}}
基准工况（$\rho=20\%$）下的计算结果如表 \ref{tab:safe_payloads} 与图 \ref{fig:safe_payloads_heatmap} 所示。

\begin{table}[H]
\centering
\small
\caption{各服务区航路属性与三机型最大安全载荷（$\rho=20\%$）}
\label{tab:safe_payloads}
\begin{tabularx}{\textwidth}{l c c c c c >{\centering\arraybackslash}X}
\toprule
\textbf{服务区} & \textbf{大圆距离} & \textbf{巡航海拔} & \textbf{A型安全载荷} & \textbf{B型安全载荷} & \textbf{C型安全载荷} & \textbf{主要制约瓶颈} \\
\textbf{编号} & ($\text{m}$) & ($\text{m}$) & ($\text{kg}$) & ($\text{kg}$) & ($\text{kg}$) & (机型 C) \\
\midrule
S001 & 3063.41 & 273.06 & 25.00 & 30.00 & 80.00 & 结构上限 \\
S002 & 7435.28 & 308.95 & 25.00 & 30.00 & \textbf{68.55} & 电池能量 \\
S003 & 7261.08 & 561.80 & 25.00 & 30.00 & \textbf{68.24} & 电池能量 \\
S004 & 7745.66 & 398.24 & 25.00 & 30.00 & \textbf{63.26} & 电池能量 \\
S005 & 5302.02 & 326.41 & 25.00 & 30.00 & 80.00 & 结构上限 \\
S006 & 3181.03 & 366.10 & 25.00 & 30.00 & 80.00 & 结构上限 \\
S007 & 4533.83 & 511.97 & 25.00 & 30.00 & 80.00 & 结构上限 \\
S008 & 8109.47 & 372.91 & 25.00 & \textbf{28.63} & \textbf{58.44} & 电池能量(B\&C) \\
S009 & 5722.68 & 293.23 & 25.00 & 30.00 & 80.00 & 结构上限 \\
S010 & 4778.44 & 435.64 & 25.00 & 30.00 & 80.00 & 结构上限 \\
S011 & 2817.15 & 295.36 & 25.00 & 30.00 & 80.00 & 结构上限 \\
S012 & 7398.45 & 413.30 & 25.00 & 30.00 & \textbf{68.01} & 电池能量 \\
S013 & 4848.71 & 399.18 & 25.00 & 30.00 & 80.00 & 结构上限 \\
S014 & 5410.92 & 591.95 & 25.00 & 30.00 & 80.00 & 结构上限 \\
S015 & 6019.11 & 571.80 & 25.00 & 30.00 & 80.00 & 结构上限 \\
\bottomrule
\end{tabularx}
\end{table}

\begin{figure}[H]
\centering
\SafeGraphic[width=0.72\textwidth]{q1_payload_heatmap.png}
\caption{各服务区与三机型最大安全载荷矩阵}
\label{fig:safe_payloads_heatmap}
\end{figure}

\begin{figure}[H]
\centering
\SafeGraphic[width=0.85\textwidth]{q1_distance_payload.png}
\caption{航线距离与最大安全载荷的关系}
\label{fig:distance_payload}
\end{figure}

结合表 \ref{tab:safe_payloads}、图 \ref{fig:safe_payloads_heatmap} 与图 \ref{fig:distance_payload}，三种机型的瓶颈差异很明显。

\begin{enumerate}[leftmargin=*]
    \item \textbf{A 型}：15 个服务区的安全载荷全部等于结构上限 $25.00\,\text{kg}$。原因是 A 型的标定单位距离电耗最低（$0.18\,\text{Wh/m}$）：即使载满 $25\,\text{kg}$ 飞最远的 S008，往返电耗也只有 $3.399\,\text{kWh}$，低于 $4.5\times0.8=3.60\,\text{kWh}$ 的上限。也就是说 A 型的载荷上限完全由结构载重决定，能量和航程都不是瓶颈。需要指出，S008 处的能量裕度只剩 $0.2\,\text{kWh}$，已不宽裕。

    \item \textbf{B 型}：标定单位距离电耗 $0.143\,\text{Wh/m}$，在 14 个服务区都能达到 $30.00\,\text{kg}$。只有 S008 例外：载满 $30\,\text{kg}$ 时往返电耗 $3.299\,\text{kWh}$，超过 $4\times0.8=3.20\,\text{kWh}$ 的上限，二分求得的临界载荷为 $28.63\,\text{kg}$，降幅 $4.6\%$。

    \item \textbf{C 型}：额定装载 $80.00\,\text{kg}$、容积 $0.250\,\text{m}^3$，但空机含电池质量就有 $69.9\,\text{kg}$，标定单位距离电耗 $0.308\,\text{Wh/m}$，约为 B 型的 2.2 倍。在 S001、S005、S006、S007、S009、S010、S011、S013、S014、S015 这 10 个服务区，满载 $80\,\text{kg}$ 的往返电耗仍低于 $6.40\,\text{kWh}$ 的上限；而在距离超过 $7.2\,\text{km}$ 的 S002、S003、S004、S008、S012 五点，满载往返电耗在 $7.37\sim8.07\,\text{kWh}$ 之间，远超上限，安全载荷被压到 $68.55$、$68.24$、$63.26$、$58.44$、$68.01\,\text{kg}$，其中 S008 降幅达 $27.0\%$。C 型的瓶颈是电池能量，且随距离恶化得比 B 型更快。

    \item \textbf{共同特征}：三种机型在所有服务区都没有触发假设 2 的航程校核。这说明在本场景的距离尺度（单程不超过 $8.11\,\text{km}$）下，决定投送能力的是电池能量和结构载重两项。
\end{enumerate}

\subsection{5.3 货箱组批：集合划分模型与状态压缩求解}

\subsubsection{5.3.1 集合划分模型}
题面规定“每个货箱不可拆分且只能安排一次”“不得跨服务区组批”。这两条使得 15 个服务区在单点往返体制下彼此解耦，整体问题分解为 15 个独立的集合划分问题\cite{balas1976, nemhauser1988}。

对服务区 $S_i$，记其货箱子集为 $\mathcal{B}_i$（$|\mathcal{B}_i| \le 15$），$\mathcal{R}_i$ 为所有合规候选架次的集合。候选架次 $r$ 由机型 $g(r)$ 和一个非空货箱子集 $\mathcal{B}_r \subseteq \mathcal{B}_i$ 唯一确定，且须满足三条约束：

\begin{align}
\text{质量：} \quad & q_r = \sum_{b \in \mathcal{B}_r} w_b \le \min\left(Q_{g(r)}, Q_{safe, i, g(r)}\right) \\
\text{体积：} \quad & v_r = \sum_{b \in \mathcal{B}_r} v_b \le V_{g(r)} \\
\text{返航电量：} \quad & E_{ir}(q_r) \le (1 - \rho_{g(r)}) E_{g(r)}^{use}
\end{align}

令 $x_{ir} \in \{0,1\}$ 表示是否采用候选架次 $r$，则服务区 $S_i$ 的组批模型为：

\begin{align}
\min \quad & \bm{J}_i = \left( \sum_{r \in \mathcal{R}_i} x_{ir}, \; \sum_{r \in \mathcal{R}_i} E_{ir} x_{ir}, \; \sum_{r \in \mathcal{R}_i} T_{ir} x_{ir} \right) \label{eq:obj_partition} \\
\text{s.t.} \quad & \sum_{r \in \mathcal{R}_i : b \in \mathcal{B}_r} x_{ir} = 1, \quad \forall b \in \mathcal{B}_i \label{eq:exact_cover} \\
& x_{ir} \in \{0, 1\}, \quad \forall r \in \mathcal{R}_i
\end{align}

式 (\ref{eq:obj_partition}) 是三个指标组成的目标向量，式 (\ref{eq:exact_cover}) 保证每个货箱被恰好运送一次。

\subsubsection{5.3.2 候选架次生成与状态压缩动态规划}
各服务区的货箱数最多为 15（S001），子集规模上限 $2^{15}=32\,768$；其余服务区为 3--8 箱，子集规模 8--256。规模很小，可以用精确算法，不需要启发式方法。候选架次的生成过程是：枚举货箱子集，先用质量与体积这两个廉价条件剪枝，再对通过剪枝的子集计算实际能耗与返航 SOC 作精确判定。

在候选集上求解集合划分时，本文采用基于位掩码的状态压缩动态规划。用状态压缩求解离散子集覆盖类问题的做法可追溯到 Held 与 Karp\cite{held1962}。具体做法如下。

\begin{enumerate}[leftmargin=*]
    \item \textbf{状态表示}：用整数 $\text{mask} \in [0, 2^{|\mathcal{B}_i|}-1]$ 的二进制位表示交付状态，第 $k$ 位为 1 表示第 $k$ 个货箱已送达。
    \item \textbf{消除排列对称}：同一个架次组合会因枚举顺序不同被重复搜索。为此，每次转移只考虑包含“当前未送达货箱中编号最小者”的候选架次（$\text{first\_bit} = \text{ctz}(\sim \text{mask})$），从而把解空间压缩为每个划分只出现一次。
    \item \textbf{字典序比较}：状态转移时直接按指定的优先级元组比较三维目标向量，最终 $\text{best}[2^{|\mathcal{B}_i|}-1]$ 即为该服务区在该优先级下的精确最优解。
\end{enumerate}

\subsection{5.4 多目标权衡}

\subsubsection{5.4.1 三种代表性策略}
题面要求综合考虑往返架次数、总运输能耗和累计作业时间三个指标。这三个指标量纲不同、含义不同，也互相牵制，一般不存在单一最优解。本文不设主观权重，而是用字典序比较三种代表性策略\cite{ehrgott2005, marler2004}：

\begin{enumerate}[leftmargin=*]
    \item \textbf{策略 I 架次优先}，优先级 $(0,1,2)$：先最少化架次，再在架次相同的方案中比能耗，最后比时间。
    \item \textbf{策略 II 能耗优先}，优先级 $(1,0,2)$：先最少化总能耗。
    \item \textbf{策略 III 时间优先}，优先级 $(2,0,1)$：先最短化累计作业时间。
\end{enumerate}

\subsubsection{5.4.2 结果与取舍}
基准工况（$\rho=20\%$、$\kappa=1.0$）下，15 个服务区独立求解后合并，结果见表 \ref{tab:tradeoff}。

\begin{table}[H]
\centering
\caption{基准工况下三种策略的综合指标对比}
\label{tab:tradeoff}
\begin{tabularx}{\textwidth}{l c c c c >{\centering\arraybackslash}X}
\toprule
\textbf{优化策略} & \textbf{优化优先级} & \textbf{总往返架次} & \textbf{总运输能耗} & \textbf{累计作业时间} & \textbf{代表方案集中} \\
\textbf{名称} & (架次, 能耗, 时间) & (架次) & ($\text{kWh}$) & ($\text{s}$) & \textbf{非支配性} \\
\midrule
架次优先 & $(0, 1, 2)$ & 18 & 59.240 & 32798.5 & 是 \\
能耗优先 & $(1, 0, 2)$ & 19 & 59.142 & 34466.5 & 是 \\
时间优先 & $(2, 0, 1)$ & 18 & 59.240 & 32798.5 & 是 \\
\bottomrule
\end{tabularx}
\end{table}

\begin{figure}[H]
\centering
\SafeGraphic[width=0.72\textwidth]{q1_tradeoff.png}
\caption{三种目标优先策略在能耗与作业时间上的分布}
\label{fig:tradeoff_q1}
\end{figure}

从表 \ref{tab:tradeoff} 与图 \ref{fig:tradeoff_q1} 可以看到两点。

其一，\textbf{“时间优先”与“架次优先”给出了完全相同的解}。这是因为地面保障时间占了很大比重：基准方案的 $32798.5\,\text{s}$ 中，纯飞行时间为 $19310.5\,\text{s}$，地面作业时间为 $13488.0\,\text{s}$，占 $41.1\%$；其中每个架次固定发生的地面时间（出动准备加基础交接）合计 $8370\,\text{s}$，与架次数成正比。既然固定时间随架次数线性增长，而各服务区的货箱组合又高度离散，最少架次方案通常也就是最短时间方案。这一现象说明，“累计作业时间”在本文口径下与架次数高度相关。

其二，\textbf{“架次优先”与“能耗优先”互不支配}。两个方案的差别只出现在 S008 一个点：

\begin{enumerate}
    \item 架次优先：S008 用 1 架次 C 型运 $45\,\text{kg}$，往返电耗 $5.857\,\text{kWh}$，作业时间 $2203.8\,\text{s}$；
    \item 能耗优先：S008 拆成 2 架次 B 型，分别运 $25\,\text{kg}$ 与 $20\,\text{kg}$，往返电耗 $2.986+2.773=5.759\,\text{kWh}$，作业时间 $1965.9+1905.9=3871.8\,\text{s}$。
\end{enumerate}

也就是说，多出一个架次，总电耗反而下降 $0.098\,\text{kWh}$（$0.17\%$），代价是作业时间增加 $1668.0\,\text{s}$（约 $27.8\,\text{min}$，$5.1\%$）。拆开算并不违反直觉：C 型空机 $69.9\,\text{kg}$，单位距离电耗是 B 型的 2.2 倍；载 $45\,\text{kg}$ 后等效航程又从 $26\,\text{km}$ 降到 $20.1\,\text{km}$，单位距离电耗进一步升高。用两个较轻的 B 型架次替代一个较重的 C 型架次，虽然多飞了一整个往返，但两次往返的电耗之和仍低于原来一次。真正被牺牲的是地面作业时间——多出一架次就要多付出一次准备与交接。

考虑到 $0.098\,\text{kWh}$ 的电量节省对应约半小时的现场作业延长，在应急场景下并不划算，本文以“架次优先”的 18 架次方案作为推荐方案。

\subsection{5.5 推荐方案明细}

\subsubsection{5.5.1 18 架次分配方案}
基准工况、架次优先策略下，状态压缩动态规划给出的 18 个架次见表 \ref{tab:batches_detail}。80 个货箱全部分配，无遗漏、无重复。

\begin{table}[H]
\centering
\scriptsize
\caption{推荐方案（18 架次）各架次组批参数与安全指标}
\label{tab:batches_detail}
\begin{tabularx}{\textwidth}{l c c c c c c c >{\centering\arraybackslash}X}
\toprule
\textbf{架次} & \textbf{服务区} & \textbf{执行} & \textbf{总质量} & \textbf{总体积} & \textbf{往返飞行} & \textbf{累计作业} & \textbf{架次能耗} & \textbf{返航SOC} \\
\textbf{编号} & \textbf{编号} & \textbf{机型} & ($\text{kg}$) & ($\text{m}^3$) & \textbf{时间}($\text{s}$) & \textbf{时间}($\text{s}$) & ($\text{kWh}$) & ($\%$) \\
\midrule
Q1-001 & S001 & C & 76.0 & 0.159 & 619.4 & 1561.4 & 2.926 & 63.42 \\
Q1-002 & S001 & C & 78.0 & 0.235 & 619.4 & 1627.4 & 3.005 & 62.44 \\
Q1-003 & S002 & B & 25.0 & 0.067 & 1189.9 & 1819.9 & 2.722 & 31.95 \\
Q1-004 & S002 & C & 56.0 & 0.144 & 1235.0 & 2045.0 & 5.739 & 28.26 \\
Q1-005 & S003 & B & 25.0 & 0.067 & 1450.2 & 2080.2 & 2.779 & 30.52 \\
Q1-006 & S003 & C & 56.0 & 0.144 & 1559.8 & 2369.8 & 5.763 & 27.96 \\
Q1-007 & S004 & C & 59.0 & 0.156 & 1429.2 & 2305.2 & 6.177 & 22.79 \\
Q1-008 & S005 & C & 59.0 & 0.156 & 1004.0 & 1880.0 & 4.239 & 47.02 \\
Q1-009 & S006 & C & 59.0 & 0.156 & 684.9 & 1560.9 & 2.594 & 67.58 \\
Q1-010 & S007 & C & 45.0 & 0.129 & 1097.5 & 1907.5 & 3.410 & 57.37 \\
Q1-011 & S008 & C & 45.0 & 0.129 & 1393.8 & 2203.8 & 5.857 & 26.78 \\
Q1-012 & S009 & B & 25.0 & 0.067 & 930.8 & 1560.8 & 2.102 & 47.44 \\
Q1-013 & S010 & B & 25.0 & 0.067 & 924.6 & 1554.6 & 1.821 & 54.47 \\
Q1-014 & S011 & B & 25.0 & 0.067 & 559.6 & 1189.6 & 1.077 & 73.07 \\
Q1-015 & S012 & B & 25.0 & 0.067 & 1293.8 & 1923.8 & 2.755 & 31.12 \\
Q1-016 & S013 & B & 25.0 & 0.067 & 880.4 & 1510.4 & 1.825 & 54.39 \\
Q1-017 & S014 & B & 25.0 & 0.067 & 1238.6 & 1868.6 & 2.137 & 46.57 \\
Q1-018 & S015 & B & 25.0 & 0.067 & 1199.6 & 1829.6 & 2.311 & 42.23 \\
\midrule
\textbf{全网汇总} & \textbf{15个区} & \textbf{9C+9B} & \textbf{758.0} & \textbf{2.011} & \textbf{19310.5} & \textbf{32798.5} & \textbf{59.240} & \textbf{均值45.3\%} \\
\bottomrule
\end{tabularx}
\end{table}

\begin{figure}[H]
\centering
\SafeGraphic[width=0.95\textwidth]{q1_sortie_energy.png}
\caption{18 个架次的去程与返程能耗构成（柱顶标注机型）}
\label{fig:sortie_energy_q1}
\end{figure}

从图 \ref{fig:sortie_energy_q1} 可以看到，能耗最高的三个架次 Q1-007、Q1-011、Q1-006 全部落在 $7\,\text{km}$ 以外的服务区，且均为 C 型重载架次；而 18 个架次中有 9 个是 B 型轻载架次，单架次能耗都在 $3\,\text{kWh}$ 上下。全网能耗之所以集中在少数几个远距离重载架次上，根源在于 C 型的单位距离电耗本身就是 B 型的 2.2 倍。

\subsubsection{5.5.2 服务区汇总与电量核验}
各服务区的执行情况见表 \ref{tab:service_summary}。机型选配呈现出明显的规律。

两架次服务区中，S001 物资总量 $154\,\text{kg}$，超过 C 型单架次能力，用 2 架次 C 型完成；S002 与 S003 各 $81\,\text{kg}$，采用“1 架次 B 型（$25\,\text{kg}$）$+$ 1 架次 C 型（$56\,\text{kg}$）”的组合。之所以这样配，是因为这两个点距离都超过 $7.2\,\text{km}$，C 型单架次装不下 $81\,\text{kg}$（安全载荷只有 $68.5\,\text{kg}$ 左右），把一部分货转给 B 型反而更省。

其余均为单架次服务区：S004 至 S008 需求量在 $45\sim59\,\text{kg}$ 之间，各由 1 架次 C 型完成；S009 至 S015 需求量均为 $25\,\text{kg}$、体积 $0.067\,\text{m}^3$，各由 1 架次 B 型完成。

电量核验上，全部 18 个架次的返航 SOC 均不低于 $20\%$ 下限。最紧的是 Q1-007（S004，C 型，$59\,\text{kg}$，距离 $7.75\,\text{km}$），返航 SOC 为 $22.79\%$，相对下限只有 $2.79$ 个百分点裕度；其次是 Q1-011（S008，C 型，$45\,\text{kg}$，距离 $8.11\,\text{km}$），SOC 为 $26.78\%$。方案在逼近安全边界的同时没有越界。

\begin{table}[H]
\centering
\small
\caption{各受灾服务区任务完成情况汇总}
\label{tab:service_summary}
\begin{tabularx}{\textwidth}{l c c c c c >{\centering\arraybackslash}X}
\toprule
\textbf{服务区编号} & \textbf{货箱总数} & \textbf{总质量}($\text{kg}$) & \textbf{往返架次} & \textbf{总能耗}($\text{kWh}$) & \textbf{累计作业时间}($\text{s}$) & \textbf{未交付箱数} \\
\midrule
S001 & 15 & 154.0 & 2 & 5.931 & 3188.9 & 0 \\
S002 & 8 & 81.0 & 2 & 8.461 & 3864.9 & 0 \\
S003 & 8 & 81.0 & 2 & 8.542 & 4450.1 & 0 \\
S004 & 6 & 59.0 & 1 & 6.177 & 2305.2 & 0 \\
S005 & 6 & 59.0 & 1 & 4.239 & 1880.0 & 0 \\
S006 & 6 & 59.0 & 1 & 2.594 & 1560.9 & 0 \\
S007 & 5 & 45.0 & 1 & 3.410 & 1907.5 & 0 \\
S008 & 5 & 45.0 & 1 & 5.857 & 2203.8 & 0 \\
S009 & 3 & 25.0 & 1 & 2.102 & 1560.8 & 0 \\
S010 & 3 & 25.0 & 1 & 1.821 & 1554.6 & 0 \\
S011 & 3 & 25.0 & 1 & 1.077 & 1189.6 & 0 \\
S012 & 3 & 25.0 & 1 & 2.755 & 1923.8 & 0 \\
S013 & 3 & 25.0 & 1 & 1.825 & 1510.4 & 0 \\
S014 & 3 & 25.0 & 1 & 2.137 & 1868.6 & 0 \\
S015 & 3 & 25.0 & 1 & 2.311 & 1829.6 & 0 \\
\bottomrule
\end{tabularx}
\end{table}

\begin{figure}[H]
\centering
\SafeGraphic[width=0.92\textwidth]{q1_site_load.png}
\caption{各服务区组批方案的运输负荷}
\label{fig:site_load_q1}
\end{figure}

图 \ref{fig:site_load_q1} 把表 \ref{tab:service_summary} 的两个主要指标并排显示。左侧能耗排序与右侧架次排序并不一致：S003 与 S002 的能耗排在全部服务区的前两位（$8.54$ 与 $8.46\,\text{kWh}$），而它们各自只有 2 个架次，原因是这两点距离都在 $7.3\,\text{km}$ 上下；反观 S001，虽然承担了 15 个货箱、$154\,\text{kg}$ 的投送量，能耗却只有 $5.93\,\text{kWh}$，因为它离调度中心最近（$3.06\,\text{km}$）。可见在本场景里，能耗主要由距离而非货量决定。

\subsection{5.6 敏感性分析}

\subsubsection{\texorpdfstring{5.6.1 返航安全余量 $\rho$ 的影响}{5.6.1 返航安全余量 rho 的影响}}
返航电量下限直接决定可用能量，是本题最需要检验的参数。敏感性分析是检验模型稳健性、识别临界阈值的常规手段\cite{saltelli2008}。本文采用单因素扰动（One-At-a-Time，OAT），在 $\{10\%,15\%,20\%,25\%,30\%\}$ 上扫描 $\rho$，保持水平能耗标定系数 $\kappa=1.0$，结果见表 \ref{tab:sensitivity_reserve} 与图 \ref{fig:reserve_sensitivity}。

\begin{table}[H]
\centering
\small
\caption{返航安全余量 $\rho$ 对全局组批结果的影响}
\label{tab:sensitivity_reserve}
\begin{tabularx}{\textwidth}{c c c c c >{\raggedright\arraybackslash}X}
\toprule
\textbf{安全余量} $\rho$ & \textbf{总架次} & \textbf{总运输能耗}($\text{kWh}$) & \textbf{作业时间}($\text{s}$) & \textbf{可行状态} & \textbf{组批结构变化} \\
\midrule
$10\%$ & 18 & 59.240 & 32798.5 & 完全可行 & 与基准一致 \\
$15\%$ & 18 & 59.240 & 32798.5 & 完全可行 & 与基准一致 \\
$20\%$ (基准) & 18 & 59.240 & 32798.5 & 完全可行 & 基准解 \\
$25\%$ & 19 & 61.190 & 34592.3 & 完全可行 & S004 由 1 架次拆为 2 架次 \\
$30\%$ & 20 & 67.358 & 36616.8 & 完全可行 & S008 增为 2 架次；S002、S003 机型组合改为 C+C \\
\bottomrule
\end{tabularx}
\end{table}

\begin{figure}[H]
\centering
\SafeGraphic[width=0.8\textwidth]{q1_reserve_sensitivity.png}
\caption{返航安全余量 $\rho$ 对总架次、总能耗与作业时间的影响曲线}
\label{fig:reserve_sensitivity}
\end{figure}

由表 \ref{tab:sensitivity_reserve} 可得到三点结论。

\begin{enumerate}[leftmargin=*]
    \item \textbf{$\rho \in [10\%,20\%]$：方案不变。} 这一区间内安全载荷确有放宽，例如最紧张的 S008 处 C 型安全载荷从 $58.44\,\text{kg}$（$\rho=20\%$）提高到 $71.11\,\text{kg}$（$\rho=10\%$），但货箱规格固定，放宽后的载荷并没有让任何一个服务区减少架次，最优组批结构始终是 18 架次。方案对这一区间的余量取值不敏感。

    \item \textbf{$\rho=25\%$：S004 首次裂变。} 此时 C 型在 S004 的安全载荷降到 $55.12\,\text{kg}$，而该点货箱总重 $59.00\,\text{kg}$，超出 $3.88\,\text{kg}$，无法再用一个架次完成，必须拆成 B 型 $17\,\text{kg}$ 加 C 型 $42\,\text{kg}$ 两架次。全网架次由 18 增至 19，总能耗升至 $61.190\,\text{kWh}$，作业时间增加 $1793.8\,\text{s}$。这是本例的临界点：再紧一点，方案结构就发生变化。

    \item \textbf{$\rho=30\%$：约束全面收紧。} S008 处 C 型安全载荷被压到 $36.00\,\text{kg}$，S004 维持 2 架次；同时 S002 与 S003 的最优组合从“B$+$C”变为两架次 C 型，S008 由 1 架次 C 型拆为 C 型 $31\,\text{kg}$ 加 B 型 $14\,\text{kg}$。总架次升至 20，总能耗 $67.358\,\text{kWh}$。
\end{enumerate}

综合来看，$\rho$ 取值在 $20\%$ 以下时结果稳定，一旦超过 $20\%$ 就接连触发裂变。附件给定的 $20\%$ 下限恰好落在这个稳定区间的上沿。

\subsubsection{\texorpdfstring{5.6.2 水平能耗标定系数 $\kappa$ 的影响}{5.6.2 水平能耗标定系数 kappa 的影响}}
水平巡航能耗是本文标定的结果（假设 1），其绝对取值带有模型不确定性。为检验结论对这条假设的依赖程度，本文用系数 $\kappa \in \{0.8,1.0,1.2\}$ 对单位距离电耗整体放缩，重算全流程，结果见图 \ref{fig:energy_model_sensitivity}。

需要说明两点。第一，$\kappa$ 只放缩水平能耗率，不改变等效航程 $L_g(q)$，因此这是单因素扰动，不是严格的物理情景（若把 $\kappa$ 解释为航程的倒数变化，还应同步修改航程公式）。第二，由于爬升能耗不随 $\kappa$ 变化，总能耗并非随 $\kappa$ 严格线性变化。

\begin{itemize}[leftmargin=*]
    \item $\kappa = 0.8$：单位距离电耗下降 $20\%$，全部服务区仍可用原结构完成配送，总架次保持 18，总能耗降至 $47.903\,\text{kWh}$，作业时间不变（$32798.5\,\text{s}$）。这表明在更乐观的能耗假设下方案还有余量。
    \item $\kappa = 1.2$：单位距离电耗上升 $20\%$，可用能量被明显压缩，远距离服务区被迫增加分批，总架次升至 20，总能耗 $88.776\,\text{kWh}$，作业时间延长至 $36950.0\,\text{s}$。
\end{itemize}

两端对比说明：架次数在能耗假设偏乐观时不受影响，但在偏悲观时会增加约 $11\%$（18$\to$20）。这部分不确定性来自水平能耗标定，无法由题面数据消除，建议在现场执行时按保守取值预留余量。

\begin{figure}[H]
\centering
\SafeGraphic[width=0.8\textwidth]{q1_energy_model_sensitivity.png}
\caption{水平能耗标定系数放缩对系统整体运行指标的敏感性响应}
\label{fig:energy_model_sensitivity}
\end{figure}



\section{六、问题二：异构无人机多点多架次运输调度模型}

\subsection{6.1 问题特征与建模假设}

\subsubsection{6.1.1 与问题一的差别}
问题一里每个架次只服务一个服务区，15 个服务区互相独立，组批只需看本点的载荷、体积与电量边界。问题二取消了这个解耦条件，同时把实体机、共享电池和物资时限一并纳入，带来三处本质变化。

第一，\textbf{架次的成本与可行性取决于整条航线}。一个架次可以按 $O01 \to S_i \to S_j \to O01$ 的顺序访问多个服务区，每交付一个点，机上剩余载荷就减少一部分，后续航段的等效航程与能耗都要按新载荷重算。因此“哪些箱装在一起、按什么顺序访问”不能分开决定，一架次内部的箱组与访问序必须整体确定。这类“车辆携带无人机、无人机可独立执行投送”的路径问题在文献中称为带无人机的旅行商问题（TSP-D）\cite{agatz2018}，多趟往返的变体亦有系统研究\cite{moshref2020}；两者的共同结论是路径与装载必须联合决策。

第二，\textbf{实体机与电池是有限资源}。基地只有 8 架飞机和 14 组共享电池，同一架飞机、同一组电池在同一时刻只能服务一个架次；电池返航后还要充电才能再次投入。这使问题从“组批”变成了“组批 + 资源排程”。

第三，\textbf{物资有时限}。80 件货箱中，医疗物资的期望送达时间与首批保障货箱的截止时间按硬约束处理，其余货箱的期望送达时间用于衡量及时性。

\subsubsection{6.1.2 实体无人机与共享电池配置}
按附件《运输无人机数据.xlsx》，8 架实体机为 U01$\sim$U04（A 型）、U05$\sim$U06（B 型）、U07$\sim$U08（C 型）；共享电池共 14 组，A 型 6 组、B 型 4 组、C 型 4 组，等效完全充电时间分别为 $1800\,\text{s}$、$2400\,\text{s}$、$3000\,\text{s}$，如表 \ref{tab:resources_inventory} 所示。

\begin{table}[H]
\centering
\small
\caption{临时调度中心 $O01$ 实体运输无人机与共享电池配置}
\label{tab:resources_inventory}
\begin{tabularx}{\textwidth}{c c c c c c >{\raggedright\arraybackslash}X}
\toprule
\textbf{机型} & \textbf{实体机编号} & \textbf{机体数量} & \textbf{共享电池编号} & \textbf{电池组数} & \textbf{等效满充时间} & \textbf{载重与速度} \\
\midrule
A 型 & U01, U02, U03, U04 & 4 架 & A-B01 $\sim$ A-B06 & 6 组 & $1800\,\text{s}$（30 min） & $25\,\text{kg}$，$12\,\text{m/s}$ \\
B 型 & U05, U06 & 2 架 & B-B01 $\sim$ B-B04 & 4 组 & $2400\,\text{s}$（40 min） & $30\,\text{kg}$，$15\,\text{m/s}$ \\
C 型 & U07, U08 & 2 架 & C-B01 $\sim$ C-B04 & 4 组 & $3000\,\text{s}$（50 min） & $80\,\text{kg}$，$15\,\text{m/s}$ \\
\midrule
\textbf{合计} & -- & 8 架 & -- & 14 组 & -- & -- \\
\bottomrule
\end{tabularx}
\end{table}

这组配置并不对称。C 型载重最大（$80\,\text{kg}$），但只有 2 架飞机、4 组电池，一组亏电电池要充 50 分钟；A 型载重最小（$25\,\text{kg}$），却有 4 架飞机、6 组电池，充满只需 30 分钟。这个不对称在后面会反复起作用：把重载任务全部压在 C 型上会让基地同时缺飞机、缺电池，而把载荷拆小、交给 A/B 型反而能换来更高的并发度。电池换装与充电在网络层面的建模已有专门讨论：Poikonen 等把无人机的换电/充电处理为在车辆上瞬时完成\cite{poikonen2017}，Cokyasar 等则把自动换电站的排队与库存一并纳入混合整数模型\cite{cokyasar2021}。本文的处理介于两者之间——电池可互换但充电需要时间，因此周转节拍成为与飞行时间同等重要的约束。

\subsubsection{6.1.3 物资时限与目标次序}
逐箱清单里有两个时间字段：医疗物资有“期望送达时间”，另有 30 件货箱标注了“首批保障”并给出“首批截止时间”。按题面，医疗物资的期望送达时间与首批保障货箱的截止时间都按硬约束处理，其余货箱的期望送达时间作为软指标。具体到本数据：

\begin{enumerate}[leftmargin=*]
    \item 医疗物资 16 件全部带期望送达时间，其中 15 件同时是首批保障箱，这类合计 31 项硬约束。
    \item 首批保障箱共 30 件，每个服务区 2 件（1 件医疗物资、1 件饮用水）。扣除已计入的 15 件医疗箱，另有 15 件饮用水箱带首批截止时间。
    \item 硬约束合计 $16+30=46$ 项。
    \item 软时限集合为全部 64 件非医疗货箱（饮用水 36 件、应急食品 19 件、生活卫生用品 9 件），按应急优先系数 $w_b$ 加权计算迟到量。
\end{enumerate}

题面要求综合考虑配送及时性、全部任务完成时间、运输能耗和架次数四项指标，但未规定优先级。本文按救援逻辑排为四级字典序：

\begin{align}
\text{第一级：} \quad & \min \quad C_{\max} = \max_{p \in \mathcal{P}} r_p \notag \\
\text{第二级：} \quad & \min \quad \sum_{b \notin \text{医疗物资}} w_b \max\left(0,\, t_b - D_b^{exp}\right) \notag \\
\text{第三级：} \quad & \min \quad E_{total} \notag \\
\text{第四级：} \quad & \min \quad |\mathcal{P}| \notag
\end{align}

式中 $r_p$ 为架次 $p$ 返抵 $O01$ 的时刻，$t_b$ 与 $D_b^{exp}$ 分别为货箱 $b$ 的实际交付时刻与期望送达时间。把完工时间放在第一位，是因为灾区救援的时效性主要由“最后一件物资什么时候送到”决定；医疗与首批保障箱的时限已经作为硬约束进入可行域，不需要再进目标函数。总能耗与总架次数排在后面，是在前两级已定的前提下再压缩运行代价。

\subsection{6.2 航段物理与资源占用模型}

\subsubsection{6.2.1 航路剖面与载荷沿途递减}
设 $\mathcal{P}$ 为全部执行架次集合。架次 $p$ 的飞行路线表示为有序节点序列
\begin{equation}
\mathcal{V}_p = (v_0, v_1, \dots, v_{K_p}, v_{K_p+1})
\end{equation}
其中起点与终点均为基地 $v_0 = v_{K_p+1} = O01$，$v_k \in \mathcal{S}$（$k = 1, \dots, K_p$）为经停的服务区，$K_p \ge 1$ 为该架次访问的服务区数。

架次 $p$ 承运的货箱集合记为 $\mathcal{B}_p \subseteq \mathcal{B}$，$\mathcal{B}_{p, v_k}$ 为其中在服务区 $v_k$ 交付的部分。无人机自 $O01$ 起飞时的初始载荷为 $q_p = \sum_{b \in \mathcal{B}_p} w_b$。沿航段 $v_k \to v_{k+1}$ 巡航时，机上剩余载荷为

\begin{equation}
q_{pk} = q_p - \sum_{l=1}^k \sum_{b \in \mathcal{B}_{p, v_l}} w_b = \sum_{l=k+1}^{K_p} \sum_{b \in \mathcal{B}_{p, v_l}} w_b
\end{equation}

返程空载航段 $v_{K_p} \to O01$ 上 $q_{p, K_p} = 0$。

巡航海拔的确定需要一个高程采样规则。若沿航线按固定步长等间距取点，当航线恰好斜切一处山脊或只擦过某个像元的角点时，有可能漏掉该像元，从而低估沿线最高地面高程，使巡航海拔定得偏低。为此，本文把航段视为经纬度平面上的线段，求出它与 DEM 栅格网格的全部交点，再按交点与区间中点定位线段经过的每一个像元（像元按闭矩形处理，压边与过角都算经过）。设经过的像元集合为 $\mathcal{C}_{v_k, v_{k+1}}$，则

\begin{equation}
H_{v_k, v_{k+1}}^{\max} = \max_{c \in \mathcal{C}_{v_k, v_{k+1}}} \text{DEM}(c), \qquad H_{pk}^{cruise} = H_{v_k, v_{k+1}}^{\max} + 50\,\text{m}
\end{equation}

这一改动只影响个别掠角航段。例如 $O01 \to S007$，等距采样得到的沿线最高地面为 $451.83\,\text{m}$（对应巡航海拔 $501.83\,\text{m}$），而闭矩形相交法捕获到航线斜切的一处山脊，沿线最高地面为 $461.97\,\text{m}$，巡航海拔相应提高至 $511.97\,\text{m}$。两种方法给出的净空都满足 $50\,\text{m}$，差别在于后者不会遗漏被航线触碰的像元。

航段两端的作业高度按题面取：$O01$ 为基地地面海拔，服务区 $v_k$ 为地面高程加 $30\,\text{m}$。连续访问多个服务区时，在 $v_k$ 完成交付后从 $H_{v_k}^{work}$ 重新爬升至下一段巡航海拔。

航段飞行时间与能耗沿用问题一的模型，只是载荷换成该航段对应的剩余载荷：

\begin{equation}
t_{pk}^{flight} = \frac{h_{pk}^+}{v_{g(p)}^\uparrow} + \frac{d_{v_k, v_{k+1}}}{v_{g(p)}^c} + \frac{h_{pk}^-}{v_{g(p)}^\downarrow}
\end{equation}

\begin{equation}
E_{pk}^{hor} = E_{g(p)}^{use} \cdot \frac{d_{v_k, v_{k+1}}}{L_{g(p)}(q_{pk})}, \qquad E_{pk}^{up} = \frac{(m_{g(p)}^{empty} + q_{pk}) \cdot g \cdot h_{pk}^+}{3.6 \times 10^6 \cdot \eta_{g(p)}}
\end{equation}

架次 $p$ 的飞行总耗时与总能耗为 $T_p^{flight} = \sum_{k=0}^{K_p} t_{pk}^{flight}$，$E_p = \sum_{k=0}^{K_p} (E_{pk}^{hor} + E_{pk}^{up})$；架次总能耗须满足 $E_p \le (1-\rho)E_{g(p)}^{use}$。

\subsubsection{6.2.2 电池两阶段充电时间}
架次 $p$ 返航后，电池剩余荷电状态为 $\text{SOC}_p^{ret} = 1 - E_p / E_{g(p)}^{use}$。按附录 2，电池再次投入任务前必须充至 $100\%$，充电时间按两阶段折算：$0\sim90\%$ 阶段占等效完全充电时间 $T_{full,g}$ 的 $65\%$，$90\%\sim100\%$ 阶段占 $35\%$，各段按 SOC 增量线性折算。由此得到从 $\text{SOC}_p^{ret}$ 充满所需的时长：

\begin{equation}
c_p = \begin{cases}
T_{full, g(p)} \left( 0.65 \times \frac{0.90 - \text{SOC}_p^{ret}}{0.90} + 0.35 \right), & \text{若 } \text{SOC}_p^{ret} < 0.90 \\[0.4em]
T_{full, g(p)} \times 0.35 \times \frac{1.00 - \text{SOC}_p^{ret}}{0.10}, & \text{若 } \text{SOC}_p^{ret} \ge 0.90
\end{cases}
\label{eq:charge_model}
\end{equation}

以执行 S001 的 P00035 为例，C 型电池返航 SOC 为 $62.44\%$，代入式 (\ref{eq:charge_model}) 得 $c = 3000 \times (0.65 \times 0.2756/0.90 + 0.35) = 1647.1\,\text{s}$。充电与占用时间的非线性意味着同一组电池连续两次出动的间隔并不是常数，这一点在排程中必须显式处理。

\subsubsection{6.2.3 资源共享口径}
同一个机型的实体机之间、同一个机型的电池之间完全可互换，这一点决定了资源的建模方式。设机型 $g$ 有 $N_g$ 架实体机与 $M_g$ 组电池。架次 $p$ 对实体机的占用区间为 $[s_p, r_p)$，其中 $s_p$ 为准备开始时刻、$r_p$ 为返航时刻；对电池的占用区间为 $[s_p,\, r_p + c_p)$，即从装机起直到充至 $100\%$ 为止。由于同型实体机可互换，“任何时刻在飞的 $g$ 型架次数不超过 $N_g$”与“存在一个可行的具体指派”等价；同型电池同理。因此模型只需约束每型资源的\emph{并发占用数}，具体编号可以在求解之后用区间着色补出来，既省去大量对称的 0-1 变量，也不损失最优性。这一等价性是 6.3.2 节累计约束和 6.4.3 节着色算法的依据。

需要说明的是，本文假定同型电池的充电曲线完全相同（均按式 (\ref{eq:charge_model}) 折算）。若各电池的充电特性有差异，上述等价性不再成立，程序会在建模前显式报错而不是给出可能失真的结果。

\subsection{6.3 基于路线模式库的协同排程模型}

\subsubsection{6.3.1 路线模式}
如果直接在模型里同时决定“哪些箱同批、按什么顺序访问、用哪型飞机、何时起飞”，箱组与访问序会形成组合爆炸。本文的做法是先把前者离散化：一个\textbf{路线模式}（pattern）就是一个固定的“箱组 + 访问序 + 机型”三元组
\begin{equation}
p = \bigl(\mathcal{B}_p,\; \mathcal{V}_p,\; g(p)\bigr)
\end{equation}
它必须整体满足载重、货舱容积与返航电量约束。模式一旦确定，下面这些量就全部可以用 6.2 节的物理模型一次算清：逐箱的交付时刻偏移 $\Delta_{pb}$（相对准备开始时刻）、架次时长 $d_p$、返航时刻偏移、能耗 $E_p$、以及返航后所需充电时长 $c_p$。

这样处理的好处是把非线性的航路物理与路径组合“离线”消化掉，模型里只剩下模式的选择与时刻的排定，规模可控。这种做法与列生成中的“限制主问题”思路一致：原问题的变量（列）数量巨大，先在一部分列上求解，再逐步补充\cite{barnhart1998}。代价是模型的最优性只能相对于模式库成立，这一点在 6.6 节会明确说明。为便于描述，记 $\mathcal{L}$ 为全部模式的集合，$p$ 被选用时记 $x_p = 1$。

\subsubsection{6.3.2 决策变量与约束}
模型的两类决策是“选哪些模式”和“每个模式什么时候开始”：

\begin{itemize}[leftmargin=*]
    \item $x_p \in \{0,1\}$：模式 $p$ 是否被选用；
    \item $s_p \in [0, H]$：模式 $p$ 的准备开始时刻，取整秒，$H$ 为时间域上界。
\end{itemize}

约束共五组。

\textbf{（1）货箱恰好覆盖一次。}全部 80 件货箱必须且仅被一个选中的模式承运：
\begin{equation}
\sum_{p \in \mathcal{L}:\; b \in \mathcal{B}_p} x_p = 1, \qquad \forall b \in \mathcal{B}
\end{equation}

\textbf{（2）硬时限。}医疗物资的期望送达时间与首批保障箱的截止时间都必须在可行域内满足：
\begin{equation}
x_p = 1 \;\Longrightarrow\; s_p + \lceil \Delta_{pb} \rceil \le \lfloor D_b^{hard} \rfloor, \qquad \forall p,\; \forall b \in \mathcal{B}_p \cap \mathcal{B}_{hard}
\end{equation}
时间取整方向是保守的：交付时刻向上取整、截止时间向下取整，保证整数解回代到连续时间后仍然可行。

\textbf{（3）实体机并发容量。}对每个机型 $g$，任意时刻在飞的 $g$ 型架次数不超过该型实体机数量：
\begin{equation}
\sum_{p:\; g(p) = g} \mathbf{1}\bigl[\,s_p \le t < s_p + \lceil d_p \rceil\,\bigr] \cdot x_p \;\le\; N_g, \qquad \forall g,\; \forall t
\end{equation}

\textbf{（4）电池并发容量。}同理，电池的占用要一直持续到充满：
\begin{equation}
\sum_{p:\; g(p) = g} \mathbf{1}\bigl[\,s_p \le t < s_p + \lceil d_p + c_p \rceil\,\bigr] \cdot x_p \;\le\; M_g, \qquad \forall g,\; \forall t
\end{equation}

\textbf{（5）完工时间。}全部任务完成时间取所有已执行架次返航时刻的最大值：
\begin{equation}
C_{\max} \ge s_p + \lceil d_p \rceil, \qquad \forall p \in \mathcal{L}
\end{equation}

软时限的迟到量按箱定义。对每个非医疗货箱 $b$ 引入 $L_b \ge 0$，并要求它不小于该箱实际交付时刻与期望时间之差：
\begin{equation}
x_p = 1 \;\Longrightarrow\; L_b \ge s_p + \lceil \Delta_{pb} \rceil - \lfloor D_b^{exp} \rfloor, \qquad \forall p \ni b
\end{equation}
注意这里用的是期望送达时间而不是硬截止时间，且只对非医疗货箱累加，避免硬约束箱被重复计入软目标。

\subsubsection{6.3.3 字典序求解与阶段冻结}
6.1.3 节的四级目标不能简单加权求和——三个指标的量纲分别是秒、加权秒和千瓦时，人为定权重会掩盖取舍。多目标优化中处理这类问题有加权和、ε-约束与字典序等几种常规做法\cite{ehrgott2005}。本文采用其中的字典序（分层序列）：按顺序依次最小化四个目标，每求解完一级就把该级的取值固定为上界，再进入下一级。

以 $Z_1, Z_2, Z_3, Z_4$ 记四级目标函数，求解过程为

\begin{align}
& \min\; Z_1 \quad \to \quad Z_1^{\ast} \notag \\
& \min\; Z_2 \quad \text{s.t.} \quad Z_1 \le Z_1^{\ast} \quad \to \quad Z_2^{\ast} \notag \\
& \min\; Z_3 \quad \text{s.t.} \quad Z_1 \le Z_1^{\ast},\; Z_2 \le Z_2^{\ast} \quad \to \quad Z_3^{\ast} \notag \\
& \min\; Z_4 \quad \text{s.t.} \quad Z_1 \le Z_1^{\ast},\; Z_2 \le Z_2^{\ast},\; Z_3 \le Z_3^{\ast}
\end{align}

约束一律用“不劣于当前最好值”的等式或不等式冻结，因此后一级的优化不会牺牲前一级已经取得的指标。需要强调的是，如果某一级在限时内只找到了可行解而没有证明最优，那么被冻结的是该级的\emph{当前上界}，整条链只能称为“限定时间内的字典序最优解”，不能据此声称精确最优。本文把所有阶段的状态与下界一并输出，见 6.6 节。

\subsection{6.4 求解算法}

\subsubsection{6.4.1 模式库的生成}
初始模式库由四类来源合成：三种逐区精确集合划分（分别按架次最少、时间最短、能耗最低）、构造式贪心组批、上述方案的两区合并结果，以及全部单箱模式。对每一个候选箱组与访问序，\emph{保留所有物理可行的机型}，而不是只留其中最好的一个——因为机型选择与后续的资源排程互相牵制，过早淘汰会切断可行解的来源。

筛选只做两件事：去掉与已有模式完全重复的条目，去掉物理不可行或最早硬时限已无法满足的条目。特别地，\textbf{不按运输能耗做跨模式支配剪枝}。这样做会让模式库偏大，但能保证不会因为一个想当然的支配关系丢掉了在资源排程阶段反而更优的模式。第二轮迭代后模式库从 337 条扩到 457 条，正说明初始库并不完备。

\subsubsection{6.4.2 CP-SAT 求解}
上述模型是一组带可选区间的整数线性约束，本文用 Google OR-Tools 的 CP-SAT 求解\cite{laborie2018}。每个模式对应一个由 $x_p$ 控制的可选区间：实体机占用区间 $[s_p,\, s_p + \lceil d_p \rceil)$、电池占用区间 $[s_p,\, s_p + \lceil d_p + c_p \rceil)$，同型资源的全部区间交给累计约束（\texttt{AddCumulative}）统一传播。硬时限用条件约束只在 $x_p = 1$ 时生效；货箱覆盖写成等式和。

四级目标按 6.3.3 节依次求解。程序给每级分配不同的时间预算，四级的大致比例为 $0.55$、$0.20$、$0.20$ 与 $0.05$，并在进入下一级前把上一级的最优解作为提示（hint）注入，让后续求解从一个已知良好的起点出发。由于求解器是多线程限时运行的，同一输入在不同机器上可能得到不同的可行解，程序把每一级的状态与下界都写进输出，便于复核。

\subsubsection{6.4.3 具体资源编号的分配}
模型只约束了每型资源的并发数，还没有回答“哪架飞机、哪组电池执行哪个架次”。这一步由区间着色完成：把同型架次的占用区间按开始时刻排序，依次为每个区间分配当前最早空闲的编号（首次适配）。对区间图而言，首次适配着色所用颜色数恰好等于最大团大小，也就恰好等于累计约束给出的并发峰值，因此\textbf{着色一定能成功，而且不会额外增加资源}。这既省去了指派变量带来的大量对称，也让 6.2.3 节的“并发约束等价于可行指派”落到可执行层面。

\subsubsection{6.4.4 反馈迭代与模式邻域扩展}
初始模式库只覆盖有限的箱组与访问序。为逐步补足，本文采用反馈迭代：每一轮求解完成后，从通过审计的方案出发生成邻域候选——把多服务区架次按服务区拆开，以及在两个已选箱组之间转移一个货箱——再对每个新箱组枚举访问序并保留可行机型。每轮默认最多接纳 80 条提案（被拒绝的提案不占名额），这个上限是显式的搜索限制，不是支配性论断。带无人机的路径问题规模稍大就需要这类邻域搜索，Schermer 等的变邻域/禁忌搜索即为代表\cite{schermer2019}。

本轮第四问得到的三组划分会作为下一轮最后一级的搜索偏好（减少跨旧分组的模式），但它排在四个主目标之后，绝不会放松前三级已经冻结的上界。每轮候选都用真实的浮点回代指标评价，最终保留最好的一轮，因此后面的搜索变差不会覆盖已经拿到的好方案。

\subsection{6.5 结果与分析}

\subsubsection{6.5.1 两轮迭代的候选方案}
按上述流程运行两轮反馈，每轮各得到一个通过审计的候选方案，指标见表 \ref{tab:q2_rounds} 与图 \ref{fig:candidate_compare_q2}。

\begin{table}[H]
\centering
\small
\caption{两轮反馈迭代得到的候选方案指标对比}
\label{tab:q2_rounds}
\begin{tabularx}{\textwidth}{l c c c c >{\raggedright\arraybackslash}X}
\toprule
\textbf{候选方案} & \textbf{完工时间}($\text{s}$) & \textbf{加权迟到量} & \textbf{总能耗}($\text{kWh}$) & \textbf{架次数} & \textbf{方案说明} \\
\midrule
第 1 轮（模式库 337 条） & 6369.81 & 4391.98 & 70.227 & 27 & 有 8 件货箱迟到 \\
\textbf{第 2 轮（模式库 457 条）} & \textbf{6347.19} & \textbf{0} & 71.970 & 28 & \textbf{最终选用} \\
\bottomrule
\end{tabularx}
\end{table}

\begin{figure}[H]
\centering
\SafeGraphic[width=0.78\textwidth]{q2_candidate_compare.png}
\caption{反馈迭代候选方案在四个目标上的取值}
\label{fig:candidate_compare_q2}
\end{figure}

两轮的对比正好体现了字典序前两级之间的取舍。第 2 轮模式库扩充后，求解器用多一个架次、多 $1.74\,\text{kWh}$ 的代价，把加权迟到量从 $4391.98$ 压到 $0$，同时完工时间还略有缩短（$6369.81 \to 6347.19\,\text{s}$）。按 6.1.3 节设定的次序，完工时间优先于迟到量、迟到量优先于能耗，第 2 轮方案在第二级上严格更优，因此被选为最终方案。需要说明的是，这里并不是“多花电换准时”这么简单——第 2 轮同时也改善了第一级目标，只是改善幅度很小（$22.6\,\text{s}$）。

\subsubsection{6.5.2 最终方案明细}
最终方案共出动 28 个架次，80 件货箱全部交付。各架次的执行飞机、电池、时刻、航线、装载质量、能耗与返航电量见表 \ref{tab:sorties_detail_q2}，按准备开始时刻排序。

\begin{table}[H]
\centering
\scriptsize
\caption{问题二最终方案（28 架次）排程与资源周转明细}
\label{tab:sorties_detail_q2}
\begin{tabularx}{\textwidth}{l c c c c c c c c c >{\centering\arraybackslash}X}
\toprule
\textbf{架次} & \textbf{实体机} & \textbf{机型} & \textbf{电池} & \textbf{准备} & \textbf{起飞} & \textbf{返航} & \textbf{访问服务区} & \textbf{质量} & \textbf{能耗} & \textbf{返航SOC} \\
\textbf{编号} & \textbf{编号} & & \textbf{编号} & ($\text{s}$) & ($\text{s}$) & ($\text{s}$) & \textbf{飞行路线} & ($\text{kg}$) & ($\text{kWh}$) & ($\%$) \\
\midrule
P00005 & U08 & C & C-B02 & 0 & 450 & 2045.0 & S002 & 56 & 5.739 & 28.26 \\
P00032 & U01 & A & A-B01 & 0 & 390 & 1312.5 & S001 & 20 & 1.268 & 71.82 \\
P00035 & U07 & C & C-B01 & 0 & 540 & 1627.4 & S001 & 78 & 3.005 & 62.44 \\
P00053 & U02 & A & A-B02 & 0 & 360 & 1312.6 & S006 & 17 & 1.310 & 70.90 \\
P00095 & U03 & A & A-B03 & 0 & 360 & 1989.0 & S014 & 17 & 2.288 & 49.16 \\
P00030 & U05 & B & B-B01 & 1 & 361 & 1906.9 & S008 & 20 & 2.773 & 30.66 \\
P00146 & U04 & A & A-B04 & 1 & 331 & 1193.5 & S001 & 14 & 1.223 & 72.82 \\
P00144 & U06 & B & B-B02 & 2 & 332 & 1092.4 & S001 & 14 & 1.010 & 74.75 \\
P00003 & U06 & B & B-B03 & 1093 & 1483 & 2912.9 & S002 & 25 & 2.722 & 31.95 \\
P00140 & U04 & A & A-B05 & 1194 & 1524 & 2386.5 & S001 & 14 & 1.223 & 72.82 \\
P00071 & U01 & A & A-B06 & 1313 & 1673 & 2966.9 & S010 & 17 & 1.952 & 56.61 \\
P00118 & U07 & C & C-B03 & 1628 & 2138 & 4041.8 & S006 $\to$ S007 & 70 & 4.288 & 46.40 \\
P00020 & U05 & B & B-B04 & 1907 & 2297 & 3830.8 & S012 & 25 & 2.755 & 31.12 \\
P00137 & U02 & A & A-B04 & 2048 & 2378 & 3240.5 & S001 & 14 & 1.223 & 72.82 \\
P00044 & U08 & C & C-B04 & 2060 & 2540 & 4495.8 & S003 & 64 & 6.148 & 23.15 \\
P00089 & U03 & A & A-B01 & 2180 & 2540 & 3792.0 & S013 & 17 & 1.958 & 56.50 \\
P00057 & U04 & A & A-B02 & 2387 & 2747 & 4114.3 & S007 & 17 & 1.905 & 57.67 \\
P00026 & U06 & B & B-B02 & 2913 & 3303 & 4742.6 & S015 & 25 & 2.311 & 42.23 \\
P00107 & U01 & A & A-B03 & 3150 & 3510 & 5160.0 & S013 $\to$ S010 & 16 & 2.246 & 50.08 \\
P00049 & U02 & A & A-B05 & 3241 & 3601 & 4936.8 & S005 & 17 & 2.129 & 52.68 \\
P00028 & U05 & B & B-B01 & 3831 & 4221 & 5796.9 & S008 & 25 & 2.986 & 25.35 \\
P00041 & U03 & A & A-B06 & 4032 & 4392 & 6294.3 & S003 & 17 & 2.981 & 33.75 \\
P00009 & U07 & C & C-B01 & 4042 & 4522 & \textbf{6347.2} & S004 & 59 & 6.177 & 22.79 \\
P00098 & U04 & A & A-B04 & 4115 & 4445 & 6044.0 & S014 & 8 & 2.185 & 51.44 \\
P00052 & U08 & C & C-B02 & 4496 & 4916 & 6244.0 & S005 & 42 & 3.802 & 52.48 \\
P00014 & U06 & B & B-B03 & 4760 & 5150 & 6320.8 & S009 & 25 & 2.102 & 47.44 \\
P00077 & U02 & A & A-B01 & 4937 & 5297 & 6160.5 & S011 & 17 & 1.155 & 74.32 \\
P00080 & U01 & A & A-B02 & 5166 & 5496 & 6329.5 & S011 & 8 & 1.105 & 75.45 \\
\midrule
\textbf{全网汇总} & \textbf{8 架全部出动} & \textbf{15A+7B+6C} & \textbf{14 组全部周转} & -- & -- & \textbf{6347.2} & \textbf{含 2 条两点回路} & \textbf{758.0} & \textbf{71.970} & \textbf{均值 51.4\%} \\
\bottomrule
\end{tabularx}
\end{table}

\subsubsection{6.5.3 方案特征分析}
最终方案与问题一的 18 架次方案在结构上差别很大，值得逐条说明。

\begin{enumerate}[leftmargin=*]
    \item \textbf{架次更多但完工更快}。问题一的 18 架次方案是在单点往返体制下按“架次最少”求出的，放进本问的资源排程后并没有优势——28 个架次的方案反而把完工时间压到了 $6347.2\,\text{s}$。原因是本问的时间瓶颈不是飞行时间，而是资源的周转速度：8 架飞机同时出动的并发能力，只有在每个架次都短、且电池能快速周转时才能发挥出来。
    \item \textbf{A 型承担了最多的架次}。28 个架次里 A 型占 15 个，B 型 7 个，C 型 6 个。这与 6.1.2 节的资源结构直接对应：A 型有 4 架飞机、6 组电池、充满只要 30 分钟，最适合高频周转；C 型虽然单次能装 $80\,\text{kg}$，但只有 2 架飞机、4 组电池、充满要 50 分钟，因此只用于真正需要大载重的任务。
    \item \textbf{邻区串联没有想象中划算}。最终方案只保留了 2 条两点回路（P00118 的 $S006 \to S007$ 与 P00107 的 $S013 \to S010$），并没有把外围邻区大量串起来。原因在于串联虽然省下一个往返，却把两个服务区的交付绑在同一架次上：一旦这个架次因为资源冲突推迟，两个点都会跟着推迟。在本问“完工时间优先”的目标下，这种连带风险往往盖过了省下来的那点飞行时间。
    \item \textbf{14 组电池全部参与周转}。表 \ref{tab:sorties_detail_q2} 中 14 组电池编号全部出现，每组的循环次数在 1$\sim$3 次之间。返航 SOC 最低的是 P00009 的 $22.79\%$（S004，C 型载 $59\,\text{kg}$，距离 $7.75\,\text{km}$），最高的是 P00080 的 $75.45\%$（S011，A 型载 $8\,\text{kg}$），全部高于 $20\%$ 下限。可见方案在把电池压榨到接近极限的同时没有越过安全边界。
\end{enumerate}

\subsubsection{6.5.4 图表解读}
\begin{figure}[H]
\centering
\SafeGraphic[width=0.85\textwidth]{q2_routes.png}
\caption{最终方案的服务区访问路线网络}
\label{fig:routes_q2}
\end{figure}

图 \ref{fig:routes_q2} 给出全部 28 个架次的航线。同一节点对之间的往返画成一条双向弧线，线型与颜色区分机型。可以看到 S001（物资量 $154\,\text{kg}$）周围聚集了 6 个架次，而外围的 S003、S010、S013 等点则被两点回路顺带覆盖。

\begin{figure}[H]
\centering
\SafeGraphic[width=0.92\textwidth]{q2_drone_gantt.png}
\caption{8 架实体无人机的任务占用时间轴}
\label{fig:drone_gantt_q2}
\end{figure}

图 \ref{fig:drone_gantt_q2} 显示 8 架飞机的占用情况。8 架都在 $t \le 2\,\text{s}$ 就开始了第一个架次，初始阶段没有出现资源闲置；此后各机的架次数在 3$\sim$4 个之间，负载比较均衡。占用最满的是 U07 与 U01，都执行到 $6300\,\text{s}$ 之后，正是它们决定了全网的完工时间。

\begin{figure}[H]
\centering
\SafeGraphic[width=0.92\textwidth]{q2_battery_gantt.png}
\caption{14 组共享电池的任务占用与充电周转}
\label{fig:battery_gantt_q2}
\end{figure}

图 \ref{fig:battery_gantt_q2} 中深色段是任务占用、浅色段是充电。可以清楚看到充电段的长度并不固定：返航 SOC 越低，充电段越长。由于电池占用要一直持续到充满，电池的周转周期实际上比飞机更长，这也是为什么方案倾向于用 A 型——它的充电时间最短，同样的飞机数量能支撑更多架次。

\begin{figure}[H]
\centering
\SafeGraphic[width=0.9\textwidth]{q2_delivery.png}
\caption{80 件货箱的期望送达时间与实际交付时间}
\label{fig:delivery_q2}
\end{figure}

图 \ref{fig:delivery_q2} 以期望送达时间为横轴、实际交付时间为纵轴，虚线是按时送达的基准线。全部 80 个点都落在基准线下方，对应加权迟到量为 0；也就是说本方案不仅满足了 46 项硬时限，连软时限也全部按期送达。

\begin{figure}[H]
\centering
\SafeGraphic[width=0.9\textwidth]{q2_sortie_energy.png}
\caption{各架次的能耗与机型分布}
\label{fig:sortie_energy_q2}
\end{figure}

图 \ref{fig:sortie_energy_q2} 按能耗排序显示 28 个架次。能耗最高的是 P00009（S004，C 型载 $59\,\text{kg}$，$6.177\,\text{kWh}$），最低的是 P00144（S001，B 型载 $14\,\text{kg}$，$1.010\,\text{kWh}$），两者相差 6 倍。排序结果清楚呈现出“机型决定量级、距离决定系数”的规律：前几位清一色是 C 型远距离架次，后几位则都是 A/B 型的近程架次。

\subsection{6.6 可行性审计与最优性边界}

\subsubsection{6.6.1 独立审计}
求解完成后，方案要经过一套独立于求解器的审计（输出为 \texttt{validation.json} 与 \texttt{terrain\_clearance\_audit.csv}）：

\begin{enumerate}[leftmargin=*]
    \item \textbf{货箱覆盖}：80 件货箱唯一映射，分配率 $100\%$，无重复、无遗漏。
    \item \textbf{硬时限}：46 项硬约束合规率 $100\%$，违约数 0。
    \item \textbf{动力学回放}：逐个架次按浮点模型重算逐段载荷、时间与能耗，与求解器记录一致。
    \item \textbf{返航电量}：28 个架次的返航 SOC 在 $22.79\%\sim75.45\%$ 之间，全部高于 $20\%$ 下限。
    \item \textbf{巡航净空}：对 58 个飞行航段（32 条不同有向航段）按线段与栅格闭矩形求交重新核验，巡航海拔与沿线最高地面高程之差的最小值为 $50.000\,\text{m}$，违规航段数为 0。
    \item \textbf{资源无重叠}：8 架飞机、14 组电池的占用区间两两无交集，累计容量约束与着色结果一致。
\end{enumerate}

\subsubsection{6.6.2 最优性边界}
这一段需要把“最优”的范围说清楚。表 \ref{tab:q2_stages} 列出四级目标的求解状态：只有加权迟到量这一级被证明达到最优（下界与最优值同为 0），其余三级在限时内只找到了可行解，尚未证明最优。

\begin{table}[H]
\centering
\small
\caption{最终方案四级字典序目标的求解状态}
\label{tab:q2_stages}
\begin{tabularx}{\textwidth}{l c c c >{\raggedright\arraybackslash}X}
\toprule
\textbf{目标} & \textbf{求解状态} & \textbf{当前最优值} & \textbf{下界} & \textbf{差距说明} \\
\midrule
完工时间 & FEASIBLE & $6348\,\text{s}$ & $5990\,\text{s}$ & 尚有约 $5.6\%$ 的潜在改进空间 \\
加权迟到量 & \textbf{OPTIMAL} & 0 & 0 & 已证明最优 \\
运输能耗 & FEASIBLE & $71.970\,\text{kWh}$ & $69.258\,\text{kWh}$ & 尚有约 $3.8\%$ 的潜在改进空间 \\
架次数 & FEASIBLE & 28 & 25 & 下界为 25，未能证明是否可达 \\
\bottomrule
\end{tabularx}
\end{table}

需要明确指出三点。第一，模式库是有限的，初始库既不覆盖全部箱组，也不覆盖全部访问排列，因此上述结论只在\emph{模式库内}成立；程序始终输出 \texttt{global\_optimal = false}。第二，模型的时间与能耗都做了保守取整（时间向上、能耗放大到微千瓦时整数），这只会让可行域更紧，不会把不可行的方案判为可行，但下界也因此不能直接当作原连续问题的下界。第三，路线扩展用的是邻域启发式（拆分与单箱转移），不是带完备定价证明的列生成，所以它提升解的质量，但不构成最优性证书。

若要把结论推进到全局最优，可以在分支定价框架下做列生成：把连续的多点航线生成作为带资源约束的最短路子问题，在主问题与子问题之间迭代生成列，从而给出全局对偶下界与最优性间隙\cite{barnhart1998}。本文没有做到这一步。就应急场景而言，在数十秒内给出一个经完整可行性审计、且软硬时限全部达标的 $6347.2\,\text{s}$ 方案，已经具备实用价值。


\section{七、问题三：通信约束下的运输与中继联合调度模型（框架预留）}
{\kaishu （注：本节为全文架构预留。问题三在问题二运输调度方案基础上引入自由空间传播损耗、衰落裕量与 DEM 视距遮挡判定，协同规划中继无人机的悬停位置与服务时段。）}

\subsection{7.1 双向无线传播损耗与 DEM 视线遮挡（LOS）判定准则}
\subsection{7.2 空中中继无人机悬停选址与能耗时序耦合机理}
\subsection{7.3 连续通信保障下的时空协同排程算法}



\section{八、问题四：救援任务分区与资源配置优化模型（框架预留）}
{\kaishu （注：本节为全文架构预留。问题四以问题三固定的运输—通信联合调度方案为输入，研究 $K=2$ 与 $K=3$ 空间分区及独立执行单元的资源核算与库存缺口。）}

\subsection{8.1 原子任务超图连通性与不可拆分组划分}
\subsection{8.2 任务组独立执行下的峰值并发资源核算算法}
\subsection{\texorpdfstring{8.3 $K=2$ 与 $K=3$ 方案的资源冗余、均衡性与库存缺口评估}{8.3 K=2 与 K=3 方案的资源冗余、均衡性与库存缺口评估}}



\section{九、模型评价与推广}

\subsection{9.1 主要特点}
\begin{enumerate}[leftmargin=*]
    \item \textbf{地形与能耗模型贴合题目场景。} 巡航海拔由 DEM 沿线最高点决定而非取固定高度，爬升高度因此随航路变化；载荷通过 $3/2$ 次方关系进入能耗与航程，避免了线性续航假设在重载下低估电耗的问题。
    \item \textbf{组批问题求解到精确最优。} 利用各服务区货箱相互独立的天然解耦，用状态压缩动态规划在毫秒级得到全局最优解，不存在启发式算法的随机性和早熟风险，也便于复核。
    \item \textbf{对关键假设做了量化检验。} 水平能耗标定是本文最主要的假设，其影响通过 $\kappa$ 敏感性单独度量；结果在 $\kappa=0.8$ 时不变化，在 $\kappa=1.2$ 时架次增加 2，边界清楚。
\end{enumerate}

\subsection{9.2 局限与适用条件}
\begin{enumerate}[leftmargin=*]
    \item \textbf{水平能耗为工程标定，不是唯一真实模型。} 题面未给水平能耗公式，本文按可用能量与等效航程之比标定。$15\times3$ 的最大安全载荷与组批能耗的绝对数值都条件于这条假设，其相对大小关系（机型排序、服务区排序）则较为稳健。
    \item \textbf{“累计作业时间”是自定义口径。} 本文把它定义为准备、装载、飞行与交接时间之和，该口径下时间与架次数高度相关，导致“时间优先”与“架次优先”给出同一方案。若题目本意是墙钟完工时间，则还需引入各服务区任务能否并行、机队规模等约束，那属于问题二的范围，本文未擅自并入。
    \item \textbf{Pareto 结论限于代表解。} 表 \ref{tab:tradeoff} 中的非支配关系只在三种字典序策略给出的代表解之间成立，不等于完整 Pareto 前沿。
    \item \textbf{几何与高程采用统一近似。} 水平距离用球面大圆距离，节点作业高度取节点表海拔、未与 DEM 取值调和。在本场景的距离尺度下这些近似的影响可以忽略。
    \item \textbf{航程校核未起作用。} 假设 2 引入的单程航程校核在全部情形下均未触发，说明本场景的可行域实际由能量约束和结构载重决定。
\end{enumerate}

\subsection{9.3 问题二的额外适用条件}
\begin{enumerate}[leftmargin=*]
    \item \textbf{最优性只在有限模式库内成立。} 第二问的结论是“给定模式库内的字典序最优”，而初始模式库既不覆盖全部箱组也不覆盖全部访问排列；四级目标中只有加权迟到量一级被证明最优，其余三级在限时内只找到可行解。程序据此始终输出 \texttt{global\_optimal = false}。要给出全局最优性证书需要列生成一类的精确方法，本文没有做到。
    \item \textbf{求解器限时导致结果不可完全复现。} CP-SAT 以多线程、限时方式运行，同一输入在不同机器或不同时刻可能收敛到不同的可行解。本文把所有阶段的状态、目标值与下界一并输出，判据是审计是否通过，而不是与某次运行逐位相同。
    \item \textbf{软迟到量的统计口径。} 本文对全部 64 件非医疗货箱计算加权迟到量；若改为只统计“既非医疗也非首批保障”的货箱（49 件），数值会不同。两种口径都能自洽，但论文中必须写明用的是哪一种。
    \item \textbf{周转时间取默认值。} 附件没有给出运输机返航后的独立周转时间，本文按“机体返航即释放、下一架次可开始准备”处理，电池则必须充满才能复用；准备与装载按串行处理。若实际保障流程需要额外的检查或转场时间，完工时间会相应延长。
    \item \textbf{保守取整的影响。} 模型把时间向上取整、能耗放大到微千瓦时整数，这只会让可行域更紧，不会把不可行方案判为可行；但由此得到的下界不能直接当作原连续问题的下界。
\end{enumerate}

\subsection{9.4 对救援组织的启示}
\begin{enumerate}[leftmargin=*]
    \item \textbf{返航安全余量不宜设得过高。} 计算表明，$\rho$ 在 $20\%$ 以内方案结构稳定，一旦提高到 $25\%$ 就触发 S004 的架次裂变，提高到 $30\%$ 时总架次增加 $11\%$、总电耗增加 $13.7\%$。若现场因天气或电池低温衰减需要提高余量，应预期到架次与作业时间的明显上升，并提前准备机队与电池。
    \item \textbf{重型机与轻型机应按距离分工。} C 型单位距离电耗高，远距离作业时能量断崖明显（S008 处安全载荷比额定值低 $27\%$）；B 型空载标准航程最长、单位距离电耗最低，适合承担 $7\,\text{km}$ 以上的远距离投送。本方案中 S002、S003 采用的“B 型小载 $+$ C 型大载”组合，以及 S008 在能耗优先策略下改用两架次 B 型，都是这一分工的体现。这种多机型混合编配的思路在无人机配送问题中已有较多讨论\cite{moshref2020}。
    \item \textbf{能耗与时效的取舍量级很小。} 本例中省电与省时的对立只值 $0.098\,\text{kWh}$ 对 $1668\,\text{s}$，决策几乎不需要犹豫。但这一结论依赖于 $20\%$ 的余量基准；余量收紧后，架次增加带来的能耗与时间会同时上升，届时取舍空间会明显变小。
\end{enumerate}



\begin{thebibliography}{99}
\bibitem{caunhye2012} Caunhye A M, Nie X, Pokharel S. Optimization models in emergency logistics: A literature review[J]. Socio-Economic Planning Sciences, 2012, 46(1): 4-13.
\bibitem{holguin2012} Holguín-Veras J, Jaller M, Van Wassenhove L N, et al. On the unique features of post-disaster humanitarian logistics[J]. Journal of Operations Management, 2012, 30(7-8): 494-506.
\bibitem{murray2015} Murray C C, Chu A G. The flying sidekick traveling salesman problem: Optimization of drone-assisted parcel delivery[J]. Transportation Research Part C: Emerging Technologies, 2015, 54: 86-109.
\bibitem{otto2018} Otto A, Agatz N, Campbell J, et al. Optimization approaches for civil applications of unmanned aerial vehicles (UAVs) or aerial drones: A survey[J]. Networks, 2018, 72(4): 411-458.
\bibitem{leishman2006} Leishman J G. Principles of Helicopter Aerodynamics[M]. 2nd ed. Cambridge: Cambridge University Press, 2006: 53-112.
\bibitem{dorling2017} Dorling K, Heinrichs J, Messier G G, et al. Vehicle routing problems for drone delivery[J]. IEEE Transactions on Systems, Man, and Cybernetics: Systems, 2017, 47(1): 70-85.
\bibitem{zhang2021} Zhang J, Campbell J F, Sweeney D C, et al. Energy consumption models for delivery drones: A comparison and assessment[J]. Transportation Research Part D: Transport and Environment, 2021, 90: 102668.
\bibitem{stolaroff2018} Stolaroff J K, Samaras C, O'Neill E R, et al. Energy use and life cycle greenhouse gas emissions of drones for commercial package delivery[J]. Nature Communications, 2018, 9(1): 409.
\bibitem{liu2017} Liu Y, Sheng H, He F, et al. An energy consumption model for multi-rotor small unmanned aircraft systems[C]//2017 IEEE 20th International Conference on Intelligent Transportation Systems (ITSC). IEEE, 2017: 1-6.
\bibitem{sinnott1984} Sinnott R W. Virtues of the Haversine[J]. Sky and Telescope, 1984, 68(2): 159.
\bibitem{vincenty1975} Vincenty T. Direct and inverse solutions of geodesics on the ellipsoid with application of nested equations[J]. Survey Review, 1975, 23(176): 88-93.
\bibitem{roberge2013} Roberge V, Tarbouchi M, Labonté G. Comparison of parallel genetic algorithm and particle swarm optimization for real-time UAV path planning in complex 3D environments[J]. IEEE Transactions on Industrial Informatics, 2013, 9(1): 132-141.
\bibitem{balas1976} Balas E, Padberg M W. Set partitioning: a survey[J]. SIAM Review, 1976, 18(4): 714-760.
\bibitem{nemhauser1988} Nemhauser G L, Wolsey L A. Integer and Combinatorial Optimization[M]. New York: Wiley-Interscience, 1988.
\bibitem{held1962} Held M, Karp R M. A dynamic programming approach to sequencing problems[J]. Journal of the Society for Industrial and Applied Mathematics, 1962, 10(1): 196-210.
\bibitem{ehrgott2005} Ehrgott M. Multicriteria Optimization[M]. 2nd ed. Berlin: Springer Science \& Business Media, 2005: 115-148.
\bibitem{marler2004} Marler R T, Arora J S. Survey of multi-objective optimization methods for engineering[J]. Structural and Multidisciplinary Optimization, 2004, 26(6): 369-395.
\bibitem{saltelli2008} Saltelli A, Ratto M, Andres T, et al. Global Sensitivity Analysis: The Primer[M]. John Wiley \& Sons, 2008.
\bibitem{moshref2020} Moshref-Javadi M, Lee S, Winkenbach M. Design and evaluation of a multi-trip delivery model with truck and drones[J]. Transportation Research Part E: Logistics and Transportation Review, 2020, 136: 101887.
\bibitem{poikonen2017} Poikonen S, Wang X, Golden B. The vehicle routing problem with drones: Extended models and connections[J]. Networks, 2017, 70(1): 34-43.
\bibitem{agatz2018} Agatz N, Bouman P, Schmidt M. Optimization approaches for the traveling salesman problem with drone[J]. Transportation Science, 2018, 52(4): 965-981.
\bibitem{cokyasar2021} Cokyasar T, Dong W, Jin M, et al. Designing a drone delivery network with automated battery swapping machines[J]. Computers \& Operations Research, 2021, 129: 105177.
\bibitem{barnhart1998} Barnhart C, Johnson E L, Nemhauser G L, et al. Branch-and-price: Column generation for solving huge integer programs[J]. Operations Research, 1998, 46(3): 316-329.
\bibitem{schermer2019} Schermer D, Moeini M, Wendt O. A hybrid VNS/Tabu search algorithm for solving the vehicle routing problem with drones and en-route operations[J]. Computers \& Operations Research, 2019, 109: 134-158.
\bibitem{laborie2018} Laborie P, Rogerie J, Shaw P, et al. IBM ILOG CP optimizer for scheduling[J]. Constraints, 2018, 23(2): 210-250.
\end{thebibliography}

\end{document}
